% k8_soundness_p1.tex -- adapted copy for preprint P1 (6 November 2026).
% Source: WRITEUP/note/k8_soundness.tex, SHA-256 90bf626c8b57b7ad3ccfce6194bc08d10f8b5e1f0a981192609a417b52c5b99b
% (the note after the fix round of 7 November 2026 on its ENCBLEAN M4-M8 revision, in the note and in this copy
% alike; before that 89c54271..., the note as revised on 7 November 2026 for ENCBLEAN M4-M8; before that
% 52b04e63..., the note as revised on 6 November 2026, about 20:55–21:06 KST, after a claims review of its ENCBLEAN revision, in the
% note and in this copy alike; before that a01e82b5..., the note as revised about 20:10 KST for the Lean package
% ENCBLEAN, in the note and in this copy alike; before that revision 50bcd17a..., the note as revised after 15:40 KST for C6UPPER/noty12-2026-11-06 and then
% after a checker's round; before that round 0fe99cfe..., before the revision 248bf037... as committed
% in f2a462b67, which corrected the proof of Corollary k8s:cor:geom; the copy was first made from 197bc214... at commit
% 2054b945c). The ONLY changes are four cross-references, two footnotes, one deleted clause and two provenance edits:
% the two pointers to the companion certificate section (one naming the draft STATUS20261104/k8_section.tex) now point
% to Section~\ref{sec:cert} of P1, the trust-table row says (Section~\ref{sec:cert}) instead of (companion section), 'As
% in the rest of the note' reads 'As in Section~\ref{sec:defs}'; a footnote after 'Lean was not rerun for this section'
% points to the re-check of 6 November in Section~\ref{sec:layer}; and a footnote after 'fresh runs reproduce the
% formula files byte for byte' says what is on record for the sixfold cube files (Section~\ref{sec:enc});
% its wording was updated on 6 November 2026 after the hash-bound cake_lpr replay (C6UPPER/replay-2026-11-06), which
% recorded matching hashes for all 64 regenerated cube files (before: the 6 sampled ones only), and again after the run
% for encoder B's 32 cubes with -Y(1,2) (C6UPPER/noty12-2026-11-06), whose files it now also covers; and, after that
% replay finished, the clause ', since it changes as the replay proceeds' was deleted after 'is not repeated here'.
% Provenance edits of 6 November 2026 (after an audit of P1): the first citation of k8_soundness_checks.py in the text
% gives its path WRITEUP/note/k8_soundness_checks.py, and the source's comment below says 'in WRITEUP/note/' instead of
% 'in this directory' (edit 8); the rigor-label paragraph says that hullVerts_subset, which is not in the axiom log,
% is covered through layers_are_orbits and exists_layer (edit 9).
% Apart from these edits the copy equals the source. Since 197bc214 the mathematical text changed in the proof of
% Corollary k8s:cor:geom (f2a462b67) and in the revision for C6UPPER/noty12-2026-11-06 (Theorem k8s:thm:sound: part
% (a'), the two paragraphs after the statement and the proof; Corollary k8s:cor:geom, the cube-cover paragraph, the remarks after Lemma k8s:lem:affine and at the end of
% Section k8s:sec:reflection, the trust table and gap (G4)), in source and copy alike. In that revision the rigor tags
% of the statements and the gaps (G1)-(G3) were unchanged; (G4) gained two sentences (part (a') and the 64-cube split
% were added after the review; Lemma k8s:lem:refl and the case g rho g^-1 = rho^-1 of Lemma k8s:lem:affine are used
% only by encoder A). The ENCBLEAN revision (about 20:10 KST) changed, in source and copy alike, the rigor-label
% paragraph, the tag of Lemma k8s:lem:cube (also Lean-checked for DIMACS clause lists) and the text after its proof,
% the 64-cube paragraph, the trust table (cube-cover row, a new last row, the caption) and gaps (G1) and (G2); the edits
% 1-9 above are untouched. The fix round after the claims review of that revision (about 20:55–21:06 KST) changed, in
% source and copy alike, the source's header comment, the rigor-label paragraph, the text after the proof of Lemma
% k8s:lem:cube, the trust table's last row and gaps (G1) and (G2); the edits 1-9 are untouched. The revision of
% 7 November 2026 for ENCBLEAN M4-M8 changed, in source and copy alike, the source's header comment, the rigor-label
% paragraph, a new paragraph after the proof of Corollary k8s:cor:geom, the trust table (layer-lemma row, passage row,
% last row, caption) and gaps (G1), (G2) (with the trust items (E1)-(E6)) and (G4); the edits 1-9 are untouched.
% Its fix round (same day) changed, in source and copy alike, the source's header comment, the rigor-label paragraph
% and gap (G2) (the sentence on texts_iff_cubesUnsat64, (E2), (E4)); the edits 1-9 are untouched.
% Check with: diff WRITEUP/note/k8_soundness.tex WRITEUP/P1/k8_soundness_p1.tex
%
% k8_soundness.tex -- soundness of the SAT encodings behind the k = 8 rotational upper bounds
% (C6UPPER: sixfold, 11 layers; C7UPPER: sevenfold, 10 layers).  Draft of 6 November 2026.
% Revised 6 November 2026 (after 15:40 KST), when F_B(6,8,11) was certified unsatisfiable on its own by 64 cubes
% (C6UPPER/noty12-2026-11-06): Theorem k8s:thm:sound gains part (a'), proved without Lemma k8s:lem:refl; Corollary
% k8s:cor:geom uses (a') for encoder B; the cube-cover paragraph, the end of Section k8s:sec:reflection, a remark after
% Lemma k8s:lem:affine, the trust table and gap (G4) say that the bounds via encoder B use neither Lemma k8s:lem:refl
% nor the case g rho g^-1 = rho^-1 of Lemma k8s:lem:affine (encoder A's route, the cross-check, uses both).
% Revised 6 November 2026 (about 20:10 KST) for the Lean package ENCBLEAN (milestones M2 and M3, committed, and M9,
% not committed):
% the rigor-label paragraph says where the EncB.* theorems are and that no result here depends on them; Lemma
% k8s:lem:cube is also tagged Lean-checked for DIMACS clause lists, the only rigor tag that changed, and the text after
% its proof names EncB.not_dimacsSatisfiable_of_cover and EncB.covers_signPatterns; the 64-cube paragraph names
% EncB.encB_unsat_of_cubes64;
% the trust table's cube-cover row and a new last row (with the caption) and gaps (G1) and (G2) say what ENCBLEAN
% checks for encoder B and what it does not. The mathematical text is unchanged.
% Revised 6 November 2026 (about 20:55–21:06 KST) after a claims review of that revision: the rigor-label paragraph says
% the EncB.* items also concern DIMACS clause lists in general; the text after the proof of Lemma k8s:lem:cube calls
% it the composition of two Lean theorems; the trust table's last row says "satisfaction by the assignment of every
% combinatorial model, for even m >= 2"; (G1) says "for even m >= 2" and that 222 of the 225 grid instances agree
% byte for byte and 3 are refused where gen_b2.py raises; (G2) speaks of the Lean string the printer function computes, and no longer credits cake_lpr's
% verification for the agreement of its parser with the package's (tested on samples only). Rigor tags and the
% mathematical text are unchanged.
% Revised 7 November 2026 for ENCBLEAN M4-M8 (the Mathlib part ENCBLEAN/geo, Mathlib f897ebcf, commit 862d95ed5):
% EncB.Geo.sixfold_es8_of_texts proves the sixfold bound from the meaning of the 64 encoder-B verdicts, so for m = 6
% via encoder B the bound is Lean-checked end to end, assuming the 64 cake_lpr verdicts and the trust items E1-E6.
% Changed: the rigor-label paragraph (the two parts of ENCBLEAN; the Lean route for m = 6); a new paragraph after the
% proof of Corollary k8s:cor:geom (EncB.Geo.sixfold66_of_unsat); the trust table's layer-lemma row, passage row, last
% row and caption; gap (G1) (the step from an admissible set to an EncB.Model is now Lean-checked for even m), gap (G2)
% (now states sixfold_es8_of_texts and lists the trust items (E1)-(E6)) and the last sentence of (G4). Rigor tags of
% the statements and the mathematical text are unchanged; m = 7 and encoder A stay written. Fix round after its review
% (same day): the rigor-label paragraph (EncB.Geo items concern sixfold configurations; one cross-reference) and gap
% (G2) (what texts_iff_cubesUnsat64 shows; (E2) names the Python harnesses that wrote the hash-verdict records; (E4)
% adds a reading of the printer's IO code).
%
% An \input-able section.  It needs the macros \orient, \conv, \dirn, \lean of symmetric_extremal.tex, the theorem
% environments theorem, lemma, proposition, corollary, definition, remark, the packages amsmath, amssymb, enumitem,
% booktabs, array (for the >{...}p{...} columns of the trust table), hyperref, and the bibliography key Supp.
% Labels are prefixed k8s: (the draft k8_section.tex uses k8:).
% Compile check: k8_soundness_wrapper.tex.  The exact checks cited below are k8_soundness_checks.py (output in
% k8_soundness_checks.log), in WRITEUP/note/.
% MERGE NOTE: the section cites CHORD20261105, STATUS20261104 and WRITEUP/note as part of Supp.  The Supp entry of
% symmetric_extremal.tex (6 November) does not list them; add them to it when this section is \input there.

\providecommand{\py}[1]{{\upshape\nolinkurl{#1}}}
\providecommand{\sgn}{\operatorname{sgn}}
\providecommand{\dirv}{\operatorname{dir}}

\section{Soundness of the encodings for the \texorpdfstring{$k=8$}{k=8} rotational upper bounds}\label{k8s:sec}

The sixfold and sevenfold upper bounds for $k=8$ rest on unsatisfiability certificates for CNF formulas written by
two programs~\cite{Supp}:
\begin{itemize}
\item \emph{encoder B}, \dirn{C6UPPER/encb/gen_b2.py} (SHA-256 \texttt{3d47459a\dots}), which decides the angular
configuration of every triple by integer comparisons; at $(m,K,t)=(7,8,10)$ it reproduces byte for byte the certified
encoder-B formula of \dirn{C7UPPER} (SHA-256 \texttt{b4d29017\dots}, written by \dirn{C7UPPER/gen_b.py}; exact
computation, \dirn{C6UPPER}~\S3b);
\item \emph{encoder A}, \dirn{C7UPPER/gen.py} (SHA-256 \texttt{4564e7f0\dots}), which classifies one representative
phase vector per pattern in binary64 arithmetic.
\end{itemize}
This section proves that the formulas are \emph{sound}: every real configuration that the upper bounds exclude
yields a satisfying assignment. Encoder B is the primary object; encoder A is treated as a cross-check
(Proposition~\ref{k8s:prop:A}). The instances are $(m,K,t)=(6,8,11)$ and $(7,8,10)$; the lemmas are stated for $m\ge3$
and $4\le K\le 2m$ where this costs nothing, except that the soundness of encoder~A is proved only for
$m\in\{6,7\}$. The certificates themselves, and which of their cake\_lpr verdicts have
been replayed against hashed files, are described in Section~\ref{sec:cert}; here we prove only the implication from unsatisfiability to geometry.

\paragraph{Rigor labels.} \emph{Lean-checked} items are in \dirn{LAYERLEAN/LayerLemma.lean} (SHA-256
\texttt{afb8542d\dots}, the file as committed in \texttt{393619a57}; namespace \lean{ErdosSzekeres.LayerLemma};
Lean~4.24.0, Mathlib \texttt{f897ebcf}; the axiom check \dirn{LAYERLEAN/logs/axioms.log} reports only \lean{propext},
\lean{Classical.choice}, \lean{Quot.sound} for every cited theorem of that file except the one-line
\lean{hullVerts_subset}, which is used in \lean{layers_are_orbits} and \lean{exists_layer} and so is covered by their
checks); the declaration names below were
checked against that file, and Lean was not rerun for this section.\footnote{Note added in this paper: for this paper the file was checked again with \texttt{lake env lean} on 6 November 2026; see Section~\ref{sec:layer}, ``Build and axioms''.} Items whose Lean names begin with \lean{EncB} are in \dirn{ENCBLEAN}, a separate
package in two parts: names beginning with \lean{EncB.Geo} are in \dirn{ENCBLEAN/geo} (Lean~4.24.0, Mathlib
\texttt{f897ebcf}), the others in \dirn{ENCBLEAN/EncB} (Lean~4.24.0 with its standard library only, no Mathlib); the
axiom checks report only \lean{propext}, \lean{Classical.choice}, \lean{Quot.sound} for each of them
(\dirn{ENCBLEAN/logs/m9_d_recheck.log}; for \dirn{ENCBLEAN/geo}, pinned by \texttt{\#guard\_msgs} in
\dirn{ENCBLEAN/geo/ES8Geo/Audit.lean}, \dirn{ENCBLEAN/logs/m8_fix_audit.log}). They concern DIMACS
clause lists in general (Lemma~\ref{k8s:lem:cube}), a Lean definition of encoder~B's formula, not
\dirn{gen_b2.py} itself, and, for the names in \lean{EncB.Geo}, sixfold configurations. No result of this section depends on them. For $m=6$ they form a second route from the
certificates of encoder~B to \lean{Sixfold66} and the sixfold bound, Lean-checked end to end and independent of the
written proofs of this section, assuming the $64$ cake\_lpr verdicts and the trust items (E1)--(E6) listed under
gap~(G2) (see also~(G1)). \emph{Written proof (not refereed)} means a
complete proof given here. Most of them expand arguments reviewed on 4 November 2026 with no blocking gap found
(\dirn{C6C7REVIEW/SOUNDNESS.md}). The others have had no internal review: those marked \emph{new here} are new
arguments, and those marked \emph{not in the review} are standard facts or assembly steps that the review does not
contain. Line numbers refer to \dirn{gen_b2.py} with the SHA-256 above. Statements about what a program emits are
established by reading its source (Section~\ref{k8s:sec:trust}, gap~(G1)).

\subsection{Configurations and layers}\label{k8s:sec:setting}

Let $O$ be the origin, $e(\varphi)=(\cos\varphi,\sin\varphi)$, and $\rho$ the counterclockwise rotation by $2\pi/m$
about $O$ (the Lean \lean{rotAbout} with centre $O$ and angle $2\pi/m$). As in Section~\ref{sec:defs},
$\orient(p,q,r)=(x_q-x_p)(y_r-y_p)-(y_q-y_p)(x_r-x_p)$ (the Lean \lean{orient}). It is invariant under translations
and rotations, changes sign under reflections, and $\orient(p_{\pi(1)},p_{\pi(2)},p_{\pi(3)})=\sgn(\pi)\orient(p_1,p_2,p_3)$
for every permutation $\pi$. A finite set $X$ is in \emph{general position} if no three of its points are collinear
(Lean \lean{GeneralPosition}). It is in \emph{convex position} if it is convex independent: Lean \lean{ConvexPosition}
is defined as Mathlib's \lean{ConvexIndependent} of the inclusion of $X$, and by Mathlib's
\lean{convexIndependent_set_iff_notMem_convexHull_diff} this holds if and only if $x\notin\conv(X\setminus\{x\})$ for
every $x\in X$. A set $S$ is \emph{$K$-free} if no $K$-element subset of $S$ is in convex position, that is, if
$\neg\,$\lean{HasConvexNGon}$\;K\;S$. For a finite set $X$ let $\mathrm{hv}(X)=\{x\in X: x\notin\conv(X\setminus\{x\})\}$
(Lean \lean{hullVerts}), and let $\mathrm{peel}(X,0)=X$,
$\mathrm{peel}(X,n+1)=\mathrm{peel}(X,n)\setminus\mathrm{hv}(\mathrm{peel}(X,n))$ (Lean \lean{peel}).

\begin{definition}[nested regular $m$-gons]\label{k8s:def:nested}
A \emph{set of $t$ nested regular $m$-gons} is a set $S=L_0\cup\dots\cup L_{t-1}$, where each $L_i=\{\rho^kx_i:k\in\mathbb Z\}$
is the orbit of a point $x_i\ne O$ (a regular $m$-gon centred at $O$ with $m$ vertices), the $L_i$ are pairwise
disjoint, and $L_j\subseteq\conv(L_i)$ for $i<j$. Write $r_i=|x_i|$. General position and $K$-freeness are separate
hypotheses.
\end{definition}

\begin{lemma}[layer lemma; Lean-checked (\lean{layers_are_orbits}, \lean{card_peel}, \lean{hullVerts_eq_orbitFS_rotAbout});
the assembly below is a written step, not in the review]\label{k8s:lem:layer}
Let $m\ge1$, $K\le2m$ and $t\ge1$, and let $S$ be a finite $\rho$-invariant $K$-free set with $O\notin S$ and $|S|=mt$.
Put $L_n=\mathrm{hv}(\mathrm{peel}(S,n))$ for $n<t$. Then each $L_n$ is one $\rho$-orbit of exactly $m$ points, the $L_n$
are pairwise disjoint, $\mathrm{peel}(S,n)=L_n\cup\dots\cup L_{t-1}$ (in particular $S=L_0\cup\dots\cup L_{t-1}$), and
$L_j\subseteq\conv(L_i)$ for $i<j$. So $S$, with these layers, is a set of $t$ nested regular $m$-gons.
\end{lemma}

\begin{proof}
For every $n$ with $\mathrm{peel}(S,n)\ne\emptyset$, \lean{layers_are_orbits} gives a point $x\in\mathrm{peel}(S,n)$ (so
$x\ne O$) with $\mathrm{hv}(\mathrm{peel}(S,n))=\{\rho^kx:0\le k<m\}$, a set of exactly $m$ points, and
$\mathrm{peel}(S,n+1)\subseteq\conv\mathrm{hv}(\mathrm{peel}(S,n))$; as $\rho^m=\mathrm{id}$, this set is the orbit
$\{\rho^kx:k\in\mathbb Z\}$. Since $m(t-1)<|S|$, \lean{card_peel} gives $|\mathrm{peel}(S,n)|=|S|-mn$ for $n\le t$; so
$\mathrm{peel}(S,n)\ne\emptyset$ for $n<t$ and $\mathrm{peel}(S,t)=\emptyset$. By definition, $\mathrm{peel}(S,n)$ is the
disjoint union of $L_n$ and $\mathrm{peel}(S,n+1)$, since $\mathrm{hv}(X)\subseteq X$. Downward induction from
$\mathrm{peel}(S,t)=\emptyset$ gives $\mathrm{peel}(S,n)=L_n\sqcup\dots\sqcup L_{t-1}$, a disjoint union. For $i<j$,
$L_j\subseteq\mathrm{peel}(S,j)\subseteq\mathrm{peel}(S,i+1)\subseteq\conv(L_i)$.
\end{proof}

\begin{definition}[admissible sets]\label{k8s:def:adm}
Let $m\ge3$, $4\le K\le2m$ and $t\ge1$. A finite set $S\subseteq\mathbb R^2$ is \emph{$(m,K,t)$-admissible} if it is
$\rho$-invariant, $O\notin S$, $|S|=mt$, and $S$ is in general position and $K$-free. Its \emph{layers} are
$L_n=\mathrm{hv}(\mathrm{peel}(S,n))$, $0\le n<t$. By Lemma~\ref{k8s:lem:layer}, $S$ with these layers is a set of $t$
nested regular $m$-gons, so Definition~\ref{k8s:def:phase} below applies to it.
\end{definition}

From here on, every configuration is an admissible set with these layers. The Lean theorems are stated for
$\mathrm{hv}(\mathrm{peel}(S,n))$ and $\mathrm{peel}(S,n)$, so they apply to the layers with no translation step.

\begin{lemma}[affine images; written proof (not refereed), routine, not in the review]\label{k8s:lem:affine}
Let $g$ be an affine bijection of $\mathbb R^2$. Then $\conv(gY)=g\conv(Y)$ for every $Y\subseteq\mathbb R^2$, so
$\mathrm{hv}(gX)=g\,\mathrm{hv}(X)$ and $\mathrm{peel}(gX,n)=g\,\mathrm{peel}(X,n)$ for every finite $X$ and every $n$; and $g$
preserves general position, convex position and $K$-freeness. If moreover $g(O)=O$ and $g\rho g^{-1}\in\{\rho,\rho^{-1}\}$,
then $g$ maps every $(m,K,t)$-admissible set $S$ to an $(m,K,t)$-admissible set whose layers are
$g(L_0),\dots,g(L_{t-1})$, in this order.
\end{lemma}

\begin{proof}
$g$ and $g^{-1}$ map convex combinations to convex combinations, so $g\conv(Y)=\conv(gY)$. Hence, for $x\in X$,
$x\in\conv(X\setminus\{x\})$ if and only if $gx\in\conv(gX\setminus\{gx\})$; this gives $\mathrm{hv}(gX)=g\,\mathrm{hv}(X)$,
and, since the injective map $g$ commutes with set differences, $\mathrm{peel}(gX,n)=g\,\mathrm{peel}(X,n)$ by induction
on $n$. The same equivalence, applied to subsets, shows that $g$ preserves convex position (in the form
$x\notin\conv(X\setminus\{x\})$ for all $x$), hence $K$-freeness; $g$ and $g^{-1}$ preserve collinearity, hence general
position. If $g\rho g^{-1}=\rho^{\pm1}$, then also $g^{-1}\rho g=\rho^{\pm1}$, so $\rho(gS)=g\rho^{\pm1}(S)=gS$; here
$\rho^{-1}S=S$ because $\rho$ maps the finite set $S$ injectively into itself. Also $O=g(O)\notin gS$ and $|gS|=|S|$.
The layers of $gS$ are $\mathrm{hv}(\mathrm{peel}(gS,n))=g\,\mathrm{hv}(\mathrm{peel}(S,n))=g(L_n)$.
\end{proof}

The case $g\rho g^{-1}=\rho^{-1}$ is needed only for the reflection $\tau$ of Lemma~\ref{k8s:lem:refl}, that is, only
for parts~(a) and~(b) of Theorem~\ref{k8s:thm:sound} (formulas with the reflection unit); the bounds via encoder~B use
part~(a$'$) and do not need it. Elsewhere $g$ is a rotation about $O$, which commutes with $\rho$, or, in
Corollary~\ref{k8s:cor:geom}, a translation, for which only the first part of the lemma is used.

\begin{remark}[nested $m$-gons are admissible; written proof (not refereed), new here]\label{k8s:rem:ident}
Let $m\ge3$ and $4\le K\le2m$. If $S$ is a set of $t$ nested regular $m$-gons $L_0,\dots,L_{t-1}$
(Definition~\ref{k8s:def:nested}) that is in general position and $K$-free, then $S$ is $(m,K,t)$-admissible and
$\mathrm{hv}(\mathrm{peel}(S,n))=L_n$ for $n<t$: its layers in the sense of Definition~\ref{k8s:def:adm} are the given
ones, in the given order. This is used only to restate Theorem~\ref{k8s:thm:sound} for sets presented as nested
$m$-gons. It is not used for Corollary~\ref{k8s:cor:geom}, so it is not part of the trust base of the upper bounds.
\end{remark}

\begin{proof}
A nontrivial rotation about $O$ fixes only $O$, and $x_i\ne O$, so each orbit $L_i$ has exactly $m$ points and avoids
$O$; the $L_i$ are $\rho$-orbits and pairwise disjoint, so $S$ is $\rho$-invariant, $O\notin S$ and $|S|=mt$. Let
$U_n=L_n\cup\dots\cup L_{t-1}$; we show $\mathrm{hv}(U_n)=L_n$. Then $\mathrm{peel}(S,n)=U_n$ follows by induction on $n$,
because $U_n\setminus L_n=U_{n+1}$.
If $y\in L_j$ with $j>n$, then $y\in\conv(L_n)\subseteq\conv(U_n\setminus\{y\})$, so $y\notin\mathrm{hv}(U_n)$. Let
$x\in L_n$ and $f(y)=\langle y,x\rangle$. For $y\in L_n\setminus\{x\}$ we have $f(y)=r_n^2\cos\gamma<r_n^2=f(x)$, where
$\gamma\not\equiv0\pmod{2\pi}$ is the angle between $x$ and $y$. Hence on $\conv(L_n)$ the value $f(x)$ is attained only at
$x$ (a convex combination with value $f(x)$ puts all its weight on $x$). Every $y\in U_n\setminus\{x\}$ lies in
$\conv(L_n)$ and differs from $x$, so $f(y)<f(x)$; the maximum of the linear function $f$ on the polytope
$\conv(U_n\setminus\{x\})$ is attained at a point of $U_n\setminus\{x\}$, hence is $<f(x)$, and $x\notin\conv(U_n\setminus\{x\})$.
\end{proof}

By Lemma~\ref{k8s:lem:layer} and Remark~\ref{k8s:rem:ident}, the $(m,K,t)$-admissible sets are exactly the $K$-free
sets of $t$ nested regular $m$-gons in general position, with the same layers. Only the first direction
(Lemma~\ref{k8s:lem:layer}) enters the proofs below.

\begin{lemma}[centre exclusion; Lean-checked (\lean{centre_not_mem_of_even})]\label{k8s:lem:centre}
If $m$ is even and $S$ is $\rho$-invariant, in general position, and contains a point $p\ne O$, then $O\notin S$.
\end{lemma}

\begin{proof}
$\rho^{m/2}$ is the half-turn, so $-p\in S$, and $p,O,-p$ are three distinct collinear points.
\end{proof}

Hence a sixfold symmetric set in general position other than $\{O\}$ has $6t$ points, and $61$ points force $t\ge11$;
a sevenfold one has $7t$ or $7t+1$ points, and $65$ points force $t\ge10$. The reduction of the two upper bounds to
the propositions \lean{Sixfold66} and \lean{Sevenfold70} (every centre-free set of $66$, resp.\ $70$, points in general
position with sixfold, resp.\ sevenfold, rotational symmetry has a convex octagon) is Lean-checked:
\lean{convexNGon_of_layers}, \lean{sixfold_octagon}, \lean{sevenfold_octagon}. In the language of
Definition~\ref{k8s:def:adm}, and after translating the centre to $O$, \lean{Sixfold66} says that no
$(6,8,11)$-admissible set exists, and \lean{Sevenfold70} that no $(7,8,10)$-admissible set exists.

\begin{lemma}[strict nesting and radius order; Lean-checked (\lean{strict_nesting}, \lean{dist2_le_of_later})]\label{k8s:lem:nest}
Let $S$ be $(m,K,t)$-admissible with layers $L_n$. For $i<j$, $x\in L_i$ and $y\in L_j$ we have $\orient(x,\rho x,y)>0$;
for $i\le j$ we have $r_j\le r_i$.
\end{lemma}

\begin{proof}
$x\in\mathrm{hv}(\mathrm{peel}(S,i))$ by the definition of $L_i$, and $y\in L_j\subseteq\mathrm{peel}(S,j)$ (Lean
\lean{hullVerts_subset}). Together with $\rho$-invariance, $O\notin S$, general position, $K$-freeness, $m\ge2$ and
$K\le2m$, these are the hypotheses of \lean{strict_nesting} (layers $n=i<n'=j$) and of \lean{dist2_le_of_later}
($n=i\le n'=j$, which gives $|y|^2\le|x|^2$), both with centre $O$.
\end{proof}

\subsection{Phases, genericity and the intended assignment}\label{k8s:sec:phases}

\begin{definition}[phases, labels, genericity]\label{k8s:def:phase}
Let $S$ be a set of $t$ nested regular $m$-gons. The directions of the points of $L_i$ form a coset of
$(2\pi/m)\mathbb Z$ modulo $2\pi$, which meets $[0,2\pi/m)$ in exactly one value $\theta_i$, the \emph{phase} of $L_i$.
For $(i,k)\in\Lambda=\{0,\dots,t-1\}\times\mathbb Z/m$ put $P(i,k)=r_i\,e(\theta_i+2\pi k/m)$; then $(i,k)\mapsto P(i,k)$
is a bijection $\Lambda\to S$ and $\rho P(i,k)=P(i,k+1)$. We order $\Lambda$ lexicographically, with
$k\in\{0,\dots,m-1\}$. The configuration is \emph{normalised} if $\theta_0=0$. Let $u=\pi/m$ for odd $m$ and $u=2\pi/m$
for even $m$, and $\psi_i=\theta_i\bmod u\in[0,u)$; for odd $m$ let $h_i\in\{0,1\}$ with $\theta_i=\psi_i+h_i\pi/m$. The
configuration is \emph{generic} if $\theta_i-\theta_j\notin u\mathbb Z$ for all $i\ne j$. For a normalised configuration
this says that $\psi_1,\dots,\psi_{t-1}$ are pairwise distinct and nonzero ($\psi_0=0$); then $\theta_i\in(0,2\pi/m)$ for
$i\ge1$, and $\theta_i\ne\pi/m$ when $m$ is odd.
\end{definition}

Genericity does not depend on the order of the layers, and it is preserved by rotations about $O$.

\begin{lemma}[collinearity with the centre; written proof (not refereed), new here]\label{k8s:lem:collinear}
A set of $t$ nested regular $m$-gons is generic if and only if any two distinct points $p,q\in S$ collinear with $O$
satisfy: $m$ is even and $q=-p$ (so $p,q$ lie in one layer and $q=\rho^{m/2}p$).
\end{lemma}

\begin{proof}
$P(i,k)$ and $P(j,l)$ are collinear with $O$ if and only if their directions differ by a multiple of $\pi$, that is,
$\theta_i-\theta_j+2\pi(k-l)/m\in\pi\mathbb Z$. Now $\pi\mathbb Z+(2\pi/m)\mathbb Z=u\mathbb Z$: for odd $m$ because
$\pi=m\cdot\frac{\pi}{m}$, $\frac{2\pi}{m}=2\cdot\frac\pi m$ and $\gcd(m,2)=1$; for even $m$ because
$\pi=\frac m2\cdot\frac{2\pi}m$. So for $i\ne j$ some $P(i,k)$, $P(j,l)$ are collinear with $O$ if and only if
$\theta_i-\theta_j\in u\mathbb Z$. For $i=j$ the condition is $2(k-l)\in m\mathbb Z$: for odd $m$ this forces $k=l$, and for
even $m$ it gives $k=l$ or $l=k+m/2$, i.e.\ $P(i,l)=-P(i,k)$.
\end{proof}

\begin{lemma}[convex position is decided by orientations; written proof (not refereed), standard, not in the
review]\label{k8s:lem:convpos}
Let $X$ be finite and in general position, and $x\in X$. Then $x\in\conv(X\setminus\{x\})$ if and only if there are
distinct $a,b,c\in X\setminus\{x\}$ such that $\orient(b,c,x)$, $\orient(c,a,x)$ and $\orient(a,b,x)$ all have the sign
of $\orient(a,b,c)$. Consequently, if a bijection between two finite sets in general position preserves the sign of
$\orient$ on all triples, it maps subsets in convex position to subsets in convex position.
\end{lemma}

\begin{proof}
By Carath\'eodory's theorem in the plane, $x\in\conv(Y)$ implies $x\in\conv(T)$ for some $T\subseteq Y$ with
$|T|\le3$. If $|T|\le2$, then $x$ equals a point of $T$ (impossible, $x\notin Y$) or lies on a segment between two
points of $X$, contradicting general position. So $x\in\conv\{a,b,c\}$ with $\orient(a,b,c)\ne0$. Write
$x=\lambda_aa+\lambda_bb+\lambda_cc$ with $\lambda_a+\lambda_b+\lambda_c=1$. Since $\orient(b,c,\cdot)$ is affine and
vanishes at $b,c$, $\orient(b,c,x)=\lambda_a\orient(b,c,a)=\lambda_a\orient(a,b,c)$; likewise
$\orient(c,a,x)=\lambda_b\orient(a,b,c)$ and $\orient(a,b,x)=\lambda_c\orient(a,b,c)$. Thus $x\in\conv\{a,b,c\}$ iff all
$\lambda\ge0$, and by general position no $\lambda$ vanishes; this is the stated sign condition. Conversely, if the
sign condition holds for distinct $a,b,c\in X\setminus\{x\}$, then $\orient(a,b,c)\ne0$, so $x$ has barycentric
coordinates $\lambda_a,\lambda_b,\lambda_c$ with respect to $a,b,c$; the same identities show that they are positive,
so $x\in\conv\{a,b,c\}\subseteq\conv(X\setminus\{x\})$. For the consequence, let $f:X\to X'$ be the bijection and
$Y\subseteq X$. Subsets of a set in general position are in general position, so by the first part and the
characterisation of convex position recalled in Section~\ref{k8s:sec:setting}, $Y$ is in convex position if and only if
no $y\in Y$ and distinct $a,b,c\in Y\setminus\{y\}$ satisfy the sign condition. That condition refers only to signs of
$\orient$ on triples of $Y$, which $f$ preserves; so $Y$ is in convex position if and only if $f(Y)$ is.
\end{proof}

\begin{lemma}[genericity; written proof (not refereed)]\label{k8s:lem:generic}
Let $m\ge3$, $4\le K\le2m$.
\begin{enumerate}[label=(\alph*)]
\item If $m$ is even, every set of $t$ nested regular $m$-gons in general position is generic.
\item If $m$ is odd and there is an $(m,K,t)$-admissible set, then there is a generic, normalised
$(m,K,t)$-admissible set (with its layers as in Definition~\ref{k8s:def:adm}).
\end{enumerate}
\end{lemma}

\begin{proof}
(a) If $\theta_i-\theta_j\in u\mathbb Z$ with $i\ne j$, the proof of Lemma~\ref{k8s:lem:collinear} gives $k,l$ with
$P(i,k)$, $P(j,l)$ collinear with $O$; then $P(i,k)$, $P(j,l)$ and $P(j,l+m/2)=-P(j,l)$ are three distinct points of $S$ on
one line (the first is not in $L_j$), contradicting general position. This is the mechanism of
Lemma~\ref{k8s:lem:centre}; no perturbation is needed.

(b) Let $S$ be $(m,K,t)$-admissible, with layers $L_i$, radii $r_i$ and phases $\theta_i$ (Lemma~\ref{k8s:lem:layer},
Definition~\ref{k8s:def:phase}). For $\delta=(\delta_1,\dots,\delta_{t-1})\in\mathbb R^{t-1}$ and $\delta_0=0$ put
$P_\delta(i,k)=r_i\,e(\theta_i+\delta_i+2\pi k/m)$ and $S_\delta=\{P_\delta(\lambda):\lambda\in\Lambda\}$. For distinct labels
$\lambda,\mu,\nu$ the function $\delta\mapsto\orient(P_\delta(\lambda),P_\delta(\mu),P_\delta(\nu))$ is continuous and is
nonzero at $\delta=0$, by general position of $S$. There are finitely many label triples, so there is $\varepsilon>0$
such that for $\|\delta\|_\infty<\varepsilon$ all these orientations have their signs at $\delta=0$. Fix such a
$\delta$. Then:
(1) distinct labels give distinct points, since $P_\delta(\lambda)=P_\delta(\mu)$ would make
$\orient(P_\delta(\lambda),P_\delta(\mu),P_\delta(\nu))=0$ for any third label $\nu$; so $|S_\delta|=mt$;
(2) $S_\delta$ is in general position;
(3) $S_\delta$ is $\rho$-invariant and $O\notin S_\delta$ (each $L_i$ is rotated as a whole and $r_i>0$);
(4) by Lemma~\ref{k8s:lem:convpos} applied to $P(\lambda)\mapsto P_\delta(\lambda)$ and its inverse, $S_\delta$ is $K$-free.
So $S_\delta$ is $(m,K,t)$-admissible, and by Lemma~\ref{k8s:lem:layer} its layers are $\rho$-orbits. Its $\rho$-orbits
are the rotated layers $\{P_\delta(i,k):k\}$ (the orbits of $\rho$ on a set are determined by the set), so its layers
are these orbits in some order (in fact in the original order, because the nesting signs of Lemma~\ref{k8s:lem:nest}
are preserved, but this is not needed), and their phases are $\theta_i+\delta_i\bmod 2\pi/m$. Since genericity does not
depend on the order of the layers, $S_\delta$ is generic as soon as
\[
\delta\notin B=\bigcup_{0\le i<j\le t-1}\{\delta:\ \delta_i-\delta_j\in\theta_j-\theta_i+u\mathbb Z\}\qquad(\delta_0=0).
\]
For each pair $i<j$, the values $\delta_i-\delta_j$ on the box $(-\varepsilon,\varepsilon)^{t-1}$ lie in an interval of
length $4\varepsilon$, which meets $\theta_j-\theta_i+u\mathbb Z$ in finitely many points $c$; each set
$\{\delta_i-\delta_j=c\}$ (or $\{\delta_j=-c\}$ when $i=0$) is a proper affine hyperplane of $\mathbb R^{t-1}$. A finite
union of proper hyperplanes has empty interior, so the box contains a $\delta\notin B$. Finally rotate $S_\delta$ about
$O$ so that its outermost layer has phase $0$: by Lemma~\ref{k8s:lem:affine} (a rotation about $O$ commutes with $\rho$)
the result is admissible with the rotated layers in the same order, and since all phases change by the same constant
modulo $2\pi/m$, it is still generic. (If $t=1$ there is nothing to perturb.)
\end{proof}

For odd $m$ the perturbation is needed even though $S$ is in general position: a point and another point on the
opposite ray from $O$ are collinear with $O$, which is not a point of $S$ (the centre of an original sevenfold set is
removed in the reduction). By Lemma~\ref{k8s:lem:collinear}, after the perturbation no two points of $S$ are
collinear with $O$.

\begin{definition}[intended assignment]\label{k8s:def:assign}
Let $S$ be a generic, normalised $(m,K,t)$-admissible set, with its layers $L_n=\mathrm{hv}(\mathrm{peel}(S,n))$. Its
\emph{intended assignment} $\alpha_S$ gives
\begin{itemize}
\item the phase-order variable $Y(i,j)$, $1\le i<j\le t-1$, the value $[\psi_i<\psi_j]$;
\item for odd $m$, the variable $H(i)$, $1\le i\le t-1$, the value $[\theta_i<\pi/m]=[h_i=0]$;
\item the orientation variable $v_c$ of a triple orbit with representative $c=(c_0,c_1,c_2)$ (sorted in $\Lambda$;
Definition~\ref{k8s:def:orbits}), the
value $[\orient(P(c_0),P(c_1),P(c_2))>0]$;
\item the chain variables the values defined in Section~\ref{k8s:sec:other} (before Lemma~\ref{k8s:lem:chain}).
\end{itemize}
\end{definition}

In both encoders $\psi_0=0$ is the least phase and $h_0=0$. Encoder~B gives layer $0$ rank $0$ and $h=0$ (lines~91
and~94). In encoder~A, \py{Enc.patterns} keeps only the orders that put layer $0$ first, gives it the phase
\texttt{0.0} and no $H$ literal (\dirn{gen.py} lines 190, 211 and 218--219), and \py{ylit} returns the constants
\texttt{True}/\texttt{False} for pairs involving layer $0$ (lines 125--135); the helper \py{hlit} is never called.
Encoder~B numbers the variables
$Y(i,j)$ first (lines 36--39, lexicographically in $(i,j)$), then $H(i)$ for odd $m$ (line 40), then the orientation
variables (lines 50--55), then the chain variables on demand (lines 161--165). Encoder~A numbers the orientation
variables first, then $Y$, then $H$, then the chain variables (\py{Enc.__init__}, \py{fam_chain}).

\subsection{Triple orbits and orientation variables}\label{k8s:sec:orbits}

\begin{definition}[triple orbits, orientation variables]\label{k8s:def:orbits}
The rotation acts on $\Lambda$ by $(i,k)\mapsto(i,k+1)$ and on its $3$-element subsets; the orbits of the latter are the
\emph{triple orbits}. Each encoder has one \emph{orientation variable} $v_c$ per triple orbit, indexed by a
representative $c=(c_0,c_1,c_2)$ sorted in $\Lambda$, meant to express $\orient(P(c_0),P(c_1),P(c_2))>0$.
\end{definition}

Encoder~B's representative \py{canon}$(T)$ (lines 42--49) is the lexicographically least of the
sorted lists $\rho^{-k}T$, $(i,k)\in T$; one variable is allocated per value of \py{canon} (lines 50--55). For an ordered
triple of distinct labels, \py{olit} (lines 63--76) computes $c=$\py{canon}$(\{a,b,d\})$, takes the first point
$(i_0,k_0)$ of the list $[a,b,d]$ such that the sorted list of $\rho^{-k_0}[a,b,d]$ equals $c$, puts
$\mathit{rt}=\rho^{-k_0}[a,b,d]$ and lets $\omega$ be the permutation with $c_j=\mathit{rt}_{\omega(j)}$; it returns $v_c$ if
$\sgn\omega=+1$ and $\neg v_c$ otherwise.

\begin{lemma}[one variable per triple orbit; written proof (not refereed)]\label{k8s:lem:orbit}
\begin{enumerate}[label=(\alph*)]
\item \py{canon} is constant on each triple orbit, and \py{canon}$(T)$ lies in the orbit of $T$; so encoder~B's
orientation variables correspond bijectively to the triple orbits.
\item Under $\alpha_S$, the literal \py{olit}$(a,b,d)$ is true if and only if $\orient(P(a),P(b),P(d))>0$, for every
ordered triple of distinct labels.
\item If $\rho^s$ with $s\not\equiv0\pmod m$ maps a $3$-set $T$ onto itself, then $3\mid m$, $T=\{x,\rho^{m/3}x,\rho^{2m/3}x\}$
lies in one layer, and $\rho^s$ permutes $T$ cyclically. Consequently the sign computed in~(b) does not depend on which
admissible rotation is used.
\end{enumerate}
For $(m,t)=(6,11)$ there are $45{,}760$ triples: the $22$ triples $\{(i,k),(i,k+2),(i,k+4)\}$ form $11$ orbits of
size~$2$ and the others $7{,}623$ orbits of size~$6$, in total $7{,}634$ orbits. For $(7,10)$ all $7{,}820$ orbits of
the $54{,}740$ triples have size~$7$ (exact computation, \dirn{WRITEUP/note/k8_soundness_checks.py}; the totals agree with the review's
independent count, \dirn{C6C7REVIEW/SOUNDNESS.md}).
\end{lemma}

\begin{proof}
(a) The candidates for $\rho^rT$ are the lists $\rho^{-(k+r)}\rho^rT=\rho^{-k}T$, $(i,k)\in T$, i.e.\ the candidates for
$T$. So \py{canon}$(\rho^rT)=$\py{canon}$(T)$, and \py{canon}$(T)$ is the sorted list of some $\rho^{-k}T$. Triples in
different orbits therefore have different representatives.

(b) The loop of \py{olit} succeeds because \py{canon} is attained by one of the candidates. Let $s=-k_0$. Rotations
about $O$ preserve $\orient$ and $P(\rho^s\lambda)=\rho^sP(\lambda)$, so
\[
\orient(P(a),P(b),P(d))=\orient(P(\mathit{rt}_0),P(\mathit{rt}_1),P(\mathit{rt}_2))
=\sgn(\omega)\,\orient(P(c_0),P(c_1),P(c_2)),
\]
because $(c_0,c_1,c_2)$ is the permutation $\omega$ of $(\mathit{rt}_0,\mathit{rt}_1,\mathit{rt}_2)$ and a permutation and
its inverse have the same sign. All orientations are nonzero by general position, so \py{olit}$(a,b,d)$, which is
$v_c$ or $\neg v_c$ according to $\sgn\omega$, is true under $\alpha_S$ exactly when $\orient(P(a),P(b),P(d))>0$. The
argument holds for every admissible choice of $s$.

(c) A rotation $\rho^s\ne\mathrm{id}$ fixes no point of $\Lambda$ (it fixes only $O$), so it permutes $T$ without fixed
points, i.e.\ as a $3$-cycle. Then $\rho^{3s}$ fixes $T$ pointwise, so $\rho^{3s}=\mathrm{id}$ and $m\mid3s$ while $m\nmid s$;
hence $3\mid m$ and $s\equiv\pm m/3$. Rotations preserve layers, so $T\subseteq L_i$ for one $i$ and
$T=\{x,\rho^{m/3}x,\rho^{2m/3}x\}$. Two admissible rotations $s,s'$ for $T$ differ by a rotation fixing the set
$\rho^{s}T=c$; the two lists $\mathit{rt}$ then differ by the identity or a $3$-cycle, both even, so $\sgn\omega$ is the
same. For $m=6$ these are the triangles $\{(i,k),(i,k+2),(i,k+4)\}$, $k\in\{0,1\}$, two per layer, each orbit of size
$6/3=2$; for $m=7$ there are none. The counts follow: $(\binom{66}{3}-22)/6=7{,}623$ and $\binom{70}{3}/7=7{,}820$.
\end{proof}

Part~(c) is not needed for~(b); it shows that the literal map never identifies an orientation variable with its own
negation, which is how the stabiliser question was raised in review. Encoder~A's representative is the least sorted
list $\rho^sT$ over all $s\in\mathbb Z/m$, and its literal function \py{Enc.o} is correct by the same computation as~(b).
Its representatives coincide with encoder~B's (exact computation over all $45{,}760$ and $54{,}740$ triples, $0$
mismatches; \dirn{C6C7REVIEW/SOUNDNESS.md}); soundness does not use this.

\subsection{Directions and the angular classification}\label{k8s:sec:angular}

For $v\ne0$ let $\dirv(v)\in[0,2\pi)$ be its direction. Let $M$ be a set of layer indices, each $i\in M$ with a phase
$\theta_i\in[0,2\pi/m)$, such that the values $\psi_i=\theta_i\bmod u$, $i\in M$, are pairwise distinct; for odd $m$
define $h_i$ by $\theta_i=\psi_i+h_i\pi/m$ as in Definition~\ref{k8s:def:phase}. Only directions matter below, so when
no configuration is given we put $P(i,k)=e(\theta_i+2\pi k/m)$. The two cases used are the layers of a generic
configuration, with its points $P(i,k)$ (normalised or not), and, in Proposition~\ref{k8s:prop:A}, the layers of one
triple with the phases sampled by encoder~A. The \emph{objects} of a label $(i,k)$, $i\in M$, are the point $P(i,k)$
and its antipode $-P(i,k)$ (a direction, not necessarily a point of $S$). Put
$N=2m$ for odd $m$ and $N=m$ for even $m$, and define the cell index $J\in\{0,\dots,N-1\}$ by
\[
J(P(i,k))=(2k+h_i)\bmod 2m,\quad J(-P(i,k))=(J(P(i,k))+m)\bmod 2m\qquad(m\text{ odd}),
\]
\[
J(P(i,k))=k,\quad J(-P(i,k))=(k+m/2)\bmod m\qquad(m\text{ even}).
\]
A \emph{rank} is a map $\mathrm{rk}:M\to\mathbb Z$
with $\mathrm{rk}(i)<\mathrm{rk}(j)\iff\psi_i<\psi_j$; the \emph{key} of an object $\omega$ of layer $i\in M$ is
$\kappa(\omega)=(J(\omega),\mathrm{rk}(i))$, ordered lexicographically. These are the keys of lines 104--111.

\begin{lemma}[key order; written proof (not refereed)]\label{k8s:lem:keys}
Let $M$, the phases $\theta_i$ ($i\in M$) and a rank be as above. For objects $\omega$ of layer $i\in M$,
$\dirv(\omega)=uJ(\omega)+\psi_i$. Hence $\dirv(\omega)<\dirv(\omega')$ if and only if
$\kappa(\omega)<\kappa(\omega')$, and $\dirv(\omega)=\dirv(\omega')$ if and only if $\kappa(\omega)=\kappa(\omega')$. Two
distinct objects have equal keys if and only if $m$ is even and they are $P(i,k)$ and $-P(i,k+m/2)$.
\end{lemma}

\begin{proof}
For odd $m$, $\theta_i+2\pi k/m=\psi_i+(\pi/m)(2k+h_i)$, and reducing modulo $2\pi=2m\cdot\pi/m$ leaves
$\psi_i+(\pi/m)J$ with $J=(2k+h_i)\bmod2m$, a value in $[0,2\pi)$ since $\psi_i<\pi/m$; the antipode adds
$\pi=m\cdot\pi/m$. For even $m$, $\theta_i+2\pi k/m=\psi_i+(2\pi/m)k$ with $0\le k<m$, and the antipode adds
$\pi=(m/2)\cdot2\pi/m$. The map $(J,\psi)\mapsto uJ+\psi$ from $\{0,\dots,N-1\}\times[0,u)$ to $[0,2\pi)$ is strictly
increasing for the lexicographic order (if $J<J'$ then $uJ+\psi<u(J+1)\le uJ'+\psi'$), and replacing $\psi$ by a rank
does not change the lexicographic order. Equal keys mean equal $J$ and equal rank, i.e.\ (the $\psi_i$, $i\in M$, are
distinct) the same layer.
Within one layer, two points or two antipodes with equal $J$ have equal $k$, and a point and its own antipode have $J$ differing by $m$ or $m/2$. A point $P(i,k)$ and an antipode
$-P(i,l)$ with equal $J$ satisfy, for odd $m$, $2(k-l)\equiv m\pmod{2m}$, impossible since $m$ is odd; for even $m$,
$k\equiv l+m/2$.
\end{proof}

We call the two orientation rules \emph{rule~A} (centre inside) and \emph{rule~B} (half-plane), as in the Lean
docstrings; the letters are unrelated to the names of the encoders.

\begin{lemma}[circular dichotomy; written proof (not refereed) for (a), Lean-checked
(\lean{orient_polar}, \lean{orient_polar_centre_inside}, \lean{orient_polar_halfplane}, \lean{sin_add_lt}) for the signs
in (b) and (c)]\label{k8s:lem:dich}
Let $p_1,p_2,p_3\ne O$ be such that the six directions $\dirv(p_a)$ and $\dirv(p_a)+\pi$ are pairwise distinct modulo
$2\pi$. Read the six objects counterclockwise as a cyclic word in the letters $\mathtt P$ (point) and $\mathtt A$
(antipode).
\begin{enumerate}[label=(\alph*)]
\item The word alternates, or the three letters $\mathtt P$ are cyclically consecutive (a rotation of
$\mathtt{PPPAAA}$); the case is determined by the word.
\item (rule A) If it alternates, then $\orient(p_{\sigma(1)},p_{\sigma(2)},p_{\sigma(3)})>0$, where
$p_{\sigma(1)},p_{\sigma(2)},p_{\sigma(3)}$ are the points in counterclockwise order.
\item (rule B) If it is a rotation of $\mathtt{PPPAAA}$ and $a,b,d$ are the points in counterclockwise order from the
first $\mathtt P$ of the block, then the counterclockwise angles $x$ from $\dirv a$ to $\dirv b$ and $y$ from $\dirv b$ to
$\dirv d$ satisfy $x,y>0$ and $x+y<\pi$; if moreover $|b|\ge|a|$ and $|b|\ge|d|$, then $\orient(a,b,d)>0$.
\end{enumerate}
\end{lemma}

\begin{proof}
(a) Let $\alpha_1,\alpha_2,\alpha_3$ be the directions of the points in counterclockwise order, and $g_1,g_2,g_3>0$ the
counterclockwise gaps ($g_a$ from $\alpha_a$ to $\alpha_{a+1}$, indices modulo $3$), so $g_1+g_2+g_3=2\pi$. No gap equals
$\pi$, since then a point and the antipode of another would coincide. \emph{Case~1: all $g_a<\pi$.} Then
$g_a<\pi<g_a+g_{a+1}$ (because $g_{a+2}<\pi$), so $\alpha_a+\pi$ lies in the open arc from $\alpha_{a+1}$ to $\alpha_{a+2}$.
Each of the three open arcs between consecutive points thus contains the antipode of the opposite point, and there are
only three antipodes: the word alternates. \emph{Case~2: some gap exceeds $\pi$}; at most one can, and after a cyclic
relabelling it is $g_3$, so $g_1+g_2<\pi$. Every $\alpha_a$ lies on the closed arc of length $g_1+g_2$ starting at
$\alpha_1$, so every antipode lies on the closed arc of length $g_1+g_2$ starting at $\alpha_1+\pi$. That arc lies inside
the open arc from $\alpha_3=\alpha_1+g_1+g_2$ to $\alpha_1+2\pi$, because $\alpha_1+g_1+g_2<\alpha_1+\pi$ and
$\alpha_1+\pi+g_1+g_2<\alpha_1+2\pi$. So the word is $\mathtt{PPPAAA}$ up to rotation, the block begins at $\alpha_1$, and
$x=g_1$, $y=g_2$. The two cases give different words.

(b) Case~1. Choose representatives $\alpha<\beta<\delta<\alpha+2\pi$ of the three directions in counterclockwise order.
Then $\beta-\alpha$, $\delta-\beta$ and $\alpha-\delta+2\pi$ are the gaps, all in $(0,\pi)$, and
\lean{orient_polar_centre_inside} (with the polar form \lean{orient_polar}) gives $\orient>0$; cyclic rotation of the
arguments does not change $\orient$.

(c) Case~2 with $\alpha=\alpha_1$, $\beta=\alpha+x$, $\delta=\alpha+x+y$: \lean{orient_polar_halfplane} needs $|a|\le|b|$,
$|d|\le|b|$, $\beta-\alpha>0$, $\delta-\beta>0$ and $\delta-\alpha<\pi$; its proof uses
$\sin x+\sin y>\sin(x+y)$ (\lean{sin_add_lt}), as
\[
\orient(a,b,d)=|a||b|\sin x+|b||d|\sin y-|d||a|\sin(x+y)\ \ge\ |a||d|\bigl(\sin x+\sin y-\sin(x+y)\bigr)>0.\qedhere
\]
\end{proof}

The case in which the six directions are not distinct is excluded by hypothesis; for a generic configuration it
occurs exactly for a same-layer antipodal pair with even $m$ (Lemma~\ref{k8s:lem:keys}). The written dichotomy in the
docstring of \dirn{gen.py} omitted this third case; it is treated in Section~\ref{k8s:sec:antipodal}.

\paragraph{Patterns.} Let $c$ be an orbit representative with layer set $M_c$ and free layers $F_c=M_c\setminus\{0\}$. A
\emph{pattern} $\pi=(\mathrm{rk},h)$ consists of a bijection $\mathrm{rk}:F_c\to\{1,\dots,|F_c|\}$, extended by
$\mathrm{rk}(0)=0$, and for odd $m$ a map $h:F_c\to\{0,1\}$, extended by $h(0)=0$. Encoder~B enumerates all patterns
(lines 89--94) and forms the premise list $\mathrm{prem}(\pi)$ (lines 96--102): for $a<b$ in $F_c$ the literal $Y(a,b)$ if
$\mathrm{rk}(a)<\mathrm{rk}(b)$ and $\neg Y(a,b)$ otherwise, and for odd $m$ and $L\in F_c$ the literal $H(L)$ if $h(L)=0$
and $\neg H(L)$ otherwise. The \emph{real pattern} $\pi_S$ has $\mathrm{rk}_S(L)=1+\#\{L'\in F_c:\psi_{L'}<\psi_L\}$ and
$h_S(L)=h_L$.

\begin{lemma}[premises; written proof (not refereed)]\label{k8s:lem:prem}
All literals of $\mathrm{prem}(\pi)$ are true under $\alpha_S$ if and only if $\pi=\pi_S$. Moreover $\mathrm{rk}_S$ is a
rank on $M_c$ in the sense above.
\end{lemma}

\begin{proof}
The $\psi_L$, $L\in F_c$, are distinct and positive (genericity), so $\mathrm{rk}_S$ is a bijection onto
$\{1,\dots,|F_c|\}$ and is strictly increasing in $\psi$ also with $\mathrm{rk}_S(0)=0$, $\psi_0=0$. Under $\alpha_S$, the
pair literal for $a<b$ is true iff $[\mathrm{rk}(a)<\mathrm{rk}(b)]=[\psi_a<\psi_b]$, and the $H$ literal iff $h(L)=h_L$.
All pairwise comparisons determine a bijection onto $\{1,\dots,|F_c|\}$ uniquely.
\end{proof}

\begin{lemma}[angular clauses of encoder B; written proof (not refereed)]\label{k8s:lem:angB}
Every clause emitted at lines 82--141 of \dirn{gen_b2.py} is satisfied by $\alpha_S$.
\end{lemma}

\begin{proof}
Fix an orbit representative $c$ and its variable $v=v_c$. For each pattern $\pi$ the code computes a value
$\ell(\pi)$, either an orientation literal or \texttt{None}, and then (lines 134--141) emits the unit clause
$[\ell]$ if all patterns give the same literal $\ell\ne$\texttt{None}, and otherwise the clause
$\neg\mathrm{prem}(\pi)\vee\ell(\pi)$ for every $\pi$ with $\ell(\pi)\ne$\texttt{None}. We show that $\ell(\pi_S)$ is
\texttt{None} or true under $\alpha_S$. This suffices: $\pi_S$ is among the enumerated patterns, so a unit clause equals
$[\ell(\pi_S)]$; and a clause for $\pi\ne\pi_S$ contains the negation of a premise literal that is false under
$\alpha_S$ (Lemma~\ref{k8s:lem:prem}).

For $\pi=\pi_S$ the keys of lines 104--111 are the keys $\kappa$ of the six objects of the labels in $c$, with the rank
$\mathrm{rk}_S$ (Lemma~\ref{k8s:lem:prem}). If two keys coincide (line 113), then $m$ is even and $c$ contains a
same-layer antipodal pair (Lemma~\ref{k8s:lem:keys}), and $\ell(\pi_S)=$\texttt{None}. Otherwise the six directions are
distinct, and sorting the keys (line 118) lists the objects by increasing direction from direction $0$
(Lemma~\ref{k8s:lem:keys}); the string \py{kinds} is the cyclic word of Lemma~\ref{k8s:lem:dich}.
\begin{itemize}
\item If \py{kinds} alternates (line 121), \py{seq} $=(\sigma(1),\sigma(2),\sigma(3))$ lists the points counterclockwise,
and $\ell=v$ if $\sgn\sigma=+1$, else $\neg v$. By Lemma~\ref{k8s:lem:dich}(b),
\[
\orient(P(c_0),P(c_1),P(c_2))=\sgn(\sigma)\,\orient(P(c_{\sigma(1)}),P(c_{\sigma(2)}),P(c_{\sigma(3)}))
\]
has the sign $\sgn\sigma$, so $\ell$ is true under $\alpha_S$.
\item Otherwise the word is a rotation of $\mathtt{PPPAAA}$ (Lemma~\ref{k8s:lem:dich}(a)). In the doubled word the
substring $\mathtt{PPP}$ occurs exactly where the block starts, so \py{dbl.find('PPP')} (line 127) is the position of
its first $\mathtt P$, and $(a,b,d)$ (lines 129--130) are the points in counterclockwise order from there. If
$\mathrm{layer}(b)\le\mathrm{layer}(a),\mathrm{layer}(d)$ (line 131), Lemma~\ref{k8s:lem:nest} gives
$r_{\mathrm{layer}(b)}\ge r_{\mathrm{layer}(a)},r_{\mathrm{layer}(d)}$, Lemma~\ref{k8s:lem:dich}(c) gives
$\orient(P(a),P(b),P(d))>0$, and $\ell=$\py{olit}$(a,b,d)$ is true by Lemma~\ref{k8s:lem:orbit}(b). If the condition
fails, $\ell=$\texttt{None}.
\end{itemize}
The assertions at lines 114, 117 and 128 are not used; they held when the formulas were generated.
\end{proof}

Exhaustively over all orbits and patterns, the six integer keys are not distinct in $622$ cases (all for $m=6$),
alternate in $7{,}901$ ($m=6$) and $56{,}630$ ($m=7$) cases, form a $\mathtt{PPPAAA}$ block in $23{,}681$ and
$167{,}997$ cases, and never do anything else (exact computation, \dirn{k8_soundness_checks.py}; a check of
Lemma~\ref{k8s:lem:dich}(a), not part of the proof). For $m=6$, $t=11$ encoder~B emits $2{,}347$ angular units and
$7{,}740$ conditional angular clauses (\dirn{C6UPPER}~\S3b).

\subsection{Phase-order, nesting and chain clauses}\label{k8s:sec:other}

\begin{lemma}[phase order; written proof (not refereed)]\label{k8s:lem:order}
The clauses $\neg Y(i,j)\vee\neg Y(j,l)\vee Y(i,l)$ and $Y(i,j)\vee Y(j,l)\vee\neg Y(i,l)$, $1\le i<j<l\le t-1$ (lines
79--81), are satisfied by $\alpha_S$.
\end{lemma}

\begin{proof}
The $\psi_i$, $i\ge1$, are distinct reals, so $\neg Y(i,j)$ means $\psi_j<\psi_i$. The first clause is transitivity of
$<$. If the second failed, then $\psi_j<\psi_i$, $\psi_l<\psi_j$ and $\psi_i<\psi_l$, a cycle.
\end{proof}

These clauses say that the $Y$ variables form a strict total order; they are not used in the proofs of the other
families, which only evaluate $\alpha_S$.

\begin{lemma}[nesting units; Lean-checked (\lean{strict_nesting}) for the geometry, written proof (not refereed) for
the code]\label{k8s:lem:nestcl}
The unit clauses \py{olit}$((i,0),(i,1),(j,k))$, $0\le i<j\le t-1$, $k\in\mathbb Z/m$ (lines 143--150), are satisfied
by $\alpha_S$.
\end{lemma}

\begin{proof}
$P(i,1)=\rho P(i,0)$ and $P(j,k)\in L_j$, so Lemma~\ref{k8s:lem:nest} gives $\orient(P(i,0),P(i,1),P(j,k))>0$; apply
Lemma~\ref{k8s:lem:orbit}(b). The check at lines 148--149 only confirms that no nesting unit contradicts an angular
unit.
\end{proof}

\paragraph{Chains.} Fix a layer $i$, the pivot $p=(i,0)$ and $Q_i=\{(j,k):j\ge i\}\setminus\{p\}$ (line 159). A
\emph{$p$-chain of length $l\ge2$} is a sequence $q_1,\dots,q_l$ in $Q_i$ with $\orient(P(p),P(q_s),P(q_{s+1}))>0$ for
$1\le s<l$ and $\orient(P(q_s),P(q_{s+1}),P(q_{s+2}))>0$ for $1\le s\le l-2$. Set
$\alpha_S(F^{(i)}_l(a,b))=[$there is a $p$-chain of length $l$ with $q_1=a$, $q_2=b]$. With $L=K-1$, lines 166--190 emit,
for the pairs $(a,b)$ whose literal \py{olit}$(p,a,b)$ is not the negation of an emitted unit (\py{okpair}):
\begin{itemize}
\item (base) $F_2(a,b)\vee\neg$\py{olit}$(p,a,b)$;
\item (extension) $\neg F_l(b,d)\vee\neg$\py{olit}$(p,a,b)\vee\neg$\py{olit}$(a,b,d)\vee F_{l+1}(a,b)$ for $2\le l\le K-3$;
\item (closing) $\neg F_{K-2}(b,d)\vee\neg$\py{olit}$(p,a,b)\vee\neg$\py{olit}$(a,b,d)$;
\end{itemize}
omitting the extension and closing clauses with $a\in\{b,d\}$ or with a turn literal \py{olit}$(a,b,d)$ that is the
negation of an emitted unit. The closing clauses forbid a $p$-chain of length $K-1$, which together with $p$ would be a
convex $K$-gon.

\begin{lemma}[chain clauses; Lean-checked (\lean{chain_clauses_sound}) for the geometry, written proof (not
refereed) for the code]\label{k8s:lem:chain}
Let $m\ge3$ and $4\le K\le2m$. The chain clauses of lines 154--190 are satisfied by $\alpha_S$.
\end{lemma}

\begin{proof}
Base clause: if $\orient(P(p),P(a),P(b))>0$, then $(a,b)$ is a $p$-chain of length~$2$. Extension clause: if
$(b,d,q_3,\dots,q_l)$ is a $p$-chain and $\orient(P(p),P(a),P(b))>0$, $\orient(P(a),P(b),P(d))>0$, then
$(a,b,d,q_3,\dots,q_l)$ is a $p$-chain of length $l+1$. (Literals are evaluated by Lemma~\ref{k8s:lem:orbit}(b).)
Closing clause: otherwise there is a $p$-chain $(a,b,d,\dots)$ of length $K-1$, i.e.\ points $q_0,\dots,q_{K-2}$ of
$\mathrm{peel}(S,i)\setminus\{P(p)\}$ with $\orient(P(p),q_s,q_{s+1})>0$ and $\orient(q_s,q_{s+1},q_{s+2})>0$. Here
$P(p)\in L_i=\mathrm{hv}(\mathrm{peel}(S,i))$ by the definition of the layers; each $q_s$ is $P(\lambda)$ for a label
$\lambda\in Q_i$ of some layer $j\ge i$, so $q_s\in L_j\subseteq\mathrm{peel}(S,i)$ (Lemma~\ref{k8s:lem:layer}), and
$q_s\ne P(p)$ because $\lambda\ne p$ and $P$ is a bijection. \lean{chain_clauses_sound} (hypotheses
$m\ge3$, $4\le K\le2m$, $O\notin S$, $\rho$-invariance, general position, $K$-freeness) says that no such chain exists.
Its proof: by \lean{lin_pos_of_peel} all points of $\mathrm{peel}(S,i)$ other than $P(p)$ lie strictly on the inner side of
the tangent to the circle of radius $r_i$ at $P(p)$; by \lean{hasConvexNGon_of_fan} a pivot and $n\ge3$ chain points in an
open half-plane through the pivot satisfy $\orient>0$ for every triple of $(x_0,\dots,x_n)=(P(p),q_0,\dots,q_{n-1})$
taken in sequence order, i.e.\ $\orient(x_s,x_{s'},x_{s''})>0$ for all $s<s'<s''$ (\lean{fan_positive}; identities
\lean{orient_trans_halfplane}, \lean{fan_step}) and are therefore in convex position
(\lean{ErdosSzekeres.Suk.convexPosition_of_positive_triples}, in \dirn{ErdosSzekeres/Suk/OrientationClique.lean},
covered by the axiom check of \lean{hasConvexNGon_of_fan}); this gives $K$ points in convex position. Omitted
clauses only weaken the formula.
\end{proof}

\begin{remark}[completeness; written proof (not refereed), not used]\label{k8s:rem:complete}
Conversely, if $S$ has $K$ points $X$ in convex position, the closing clauses fail under $\alpha_S$. Let $i$ be the least
layer met by $X$ and $P(i,k)\in X$; then $\rho^{-k}X$ is in convex position, contains $P(p)$, and its other points have
labels in $Q_i$. Listing $\rho^{-k}X$ counterclockwise along its convex hull from $P(p)$ gives $P(p),q_1,\dots,q_{K-1}$,
and every three vertices of a convex polygon in general position, taken in counterclockwise order, are positively
oriented. So $(q_1,\dots,q_{K-1})$ is a $p$-chain of length $K-1$. Thus the chain family expresses $K$-freeness exactly;
soundness needs only Lemma~\ref{k8s:lem:chain}.
\end{remark}

\paragraph{Encoder A's chains.} \py{fam_chain} of \dirn{gen.py} builds the chains forwards:
$R_l(q,q')$ means a $p$-chain of length $l$ ending with $q,q'$, and it is assigned that truth value. The base,
extension and closing clauses are justified as above by appending instead of prepending. Encoder~A also uses the
emitted angular and nesting units (recorded in \py{forced}) in two ways. If the unit $[\ell]$ is emitted, it drops the
literal $\neg\ell$ from a chain clause; since $\alpha_S$ satisfies $[\ell]$ (Proposition~\ref{k8s:prop:A}, for
$m\in\{6,7\}$), it
satisfies the shortened clause exactly when it satisfies the full one. If the unit $[\neg\ell]$ is emitted, it omits
every chain clause containing $\neg\ell$, as encoder~B does; such a clause is subsumed by the unit, and omitting it is
a weakening.

\subsection{The reflection unit}\label{k8s:sec:reflection}

\begin{lemma}[reflection; written proof (not refereed)]\label{k8s:lem:refl}
Let $S$ be a generic, normalised $(m,K,t)$-admissible set, and $\tau(x,y)=(x,-y)$. Then $\tau S$ is again a generic,
normalised $(m,K,t)$-admissible set, its layers are $\tau(L_0),\dots,\tau(L_{t-1})$ in this order, and its phases are
$\theta'_0=0$ and, for $i\ge1$, $\theta'_i=2\pi/m-\theta_i$, with $\psi'_i=u-\psi_i$ and, for odd $m$, $h'_i=1-h_i$. Hence $\alpha_{\tau S}(Y(i,j))=\neg\alpha_S(Y(i,j))$ for $1\le i<j$, and for odd $m$
$\alpha_{\tau S}(H(i))=\neg\alpha_S(H(i))$. In particular exactly one of $S$ and $\tau S$ satisfies the unit $Y(1,2)$
(even $m$, $t\ge3$), resp.\ $H(1)$ (odd $m$, $t\ge2$).
\end{lemma}

\begin{proof}
$\tau$ is a linear bijection with $\tau(O)=O$ and $\tau\rho\tau^{-1}=\rho^{-1}$. By Lemma~\ref{k8s:lem:affine}, $\tau S$ is
$(m,K,t)$-admissible and its layers are $\tau(L_0),\dots,\tau(L_{t-1})$, in this order; $\tau(L_i)$ is the $\rho$-orbit of
$\tau x_i$, a regular $m$-gon of radius $r_i$. Since $\tau(r\,e(\varphi))=r\,e(-\varphi)$, the directions of
$\tau(L_i)$ are $-\theta_i+(2\pi/m)\mathbb Z$; the phase is $0$ for $i=0$ and $2\pi/m-\theta_i$ for $i\ge1$, because
$\theta_i\in(0,2\pi/m)$. Differences of phases change sign modulo $2\pi/m\in u\mathbb Z$, so genericity is kept. For
even $m$, $\psi'_i=\theta'_i=u-\psi_i$. For odd $m$, $\theta_i=\psi_i+h_i\pi/m$ with $\psi_i\in(0,\pi/m)$ gives
$\theta'_i=(\pi/m-\psi_i)+(1-h_i)\pi/m$ with $\pi/m-\psi_i\in(0,\pi/m)$, so $\psi'_i=\pi/m-\psi_i$ and $h'_i=1-h_i$. The map
$\psi\mapsto u-\psi$ reverses the order of distinct phases.
\end{proof}

The reflection need not be a symmetry of $S$: it produces another configuration of the same kind. Layer indices
must not be permuted, since they encode the radial order used by rule~B and by the nesting and chain clauses; the
reflection keeps them.

The reflection unit is used only by encoder~A, the cross-check. Encoder~B's formulas are certified unsatisfiable on
their own: at $(7,8,10)$ directly, and at $(6,8,11)$ through the $64$ cubes of Section~\ref{k8s:sec:cube}. For them
Theorem~\ref{k8s:thm:sound}(a$'$) suffices, and its proof does not use Lemma~\ref{k8s:lem:refl}.

\subsection{The soundness theorem}\label{k8s:sec:theorem}

\begin{theorem}[soundness; written proof (not refereed), resting on the Lean-checked lemmas cited above]\label{k8s:thm:sound}
Let $m\ge3$, $4\le K\le2m$ and $t\ge3$. Let $F_B$ and $F_A$ be the CNF formulas written by \dirn{gen_b2.py} and
\dirn{gen.py} on input $(m,K,t)$, and let the \emph{reflection unit} be the unit clause $Y(1,2)$ for even $m$ and
$H(1)$ for odd $m$, in each encoder's own numbering. Suppose that some $(m,K,t)$-admissible set exists, that is, a
set of $mt$ points in general position, invariant under $\rho$, not containing $O$, with no $K$ points in convex
position. By Remark~\ref{k8s:rem:ident} this holds in particular if some set of $t$ nested regular $m$-gons in general
position has no $K$ points in convex position. Then:
\begin{enumerate}[label=(\alph*)]
\item $F_B$ plus the reflection unit is satisfiable;
\item[(a$'$)] $F_B$ is satisfiable, by an argument that does not use Lemma~\ref{k8s:lem:refl};
\item if moreover $m\in\{6,7\}$, then $F_A$ plus the reflection unit is satisfiable.
\end{enumerate}
\end{theorem}

Part~(b) is restricted to $m\in\{6,7\}$ because encoder~A classifies triples in binary64 arithmetic, and the rounding
bound used in Proposition~\ref{k8s:prop:A}(i) has been computed only for these two values of $m$ (it does not depend on
$K$ or $t$). For other $m$ the same argument needs that bound for the given $m$; we make no claim about $F_A$ there.
Encoder~B uses only integer comparisons, so parts~(a) and~(a$'$) need no such bound.

Part~(a$'$) follows from~(a), since a model of $F_B$ plus a unit clause is a model of $F_B$. Its own proof below omits
the reflection step, so that the bounds via encoder~B (Corollary~\ref{k8s:cor:geom}) use neither
Lemma~\ref{k8s:lem:refl} nor the case $g\rho g^{-1}=\rho^{-1}$ of Lemma~\ref{k8s:lem:affine}. Part~(a$'$) was added on
6 November 2026, after $F_B(6,8,11)$ had been certified unsatisfiable on its own (Section~\ref{k8s:sec:cube}).

For the certified instances: at $(6,8,11)$ both formulas have $78{,}059$ variables and $2{,}387{,}083$ clauses
(SHA-256 \texttt{479e31a9\dots} for B, \texttt{542c8a3a\dots} for A), and $Y(1,2)$ is variable $1$ of $F_B$ (lines 36--39)
and variable $7{,}635=7{,}634+1$ of $F_A$; at $(7,8,10)$ both have $80{,}080$ variables and $2{,}748{,}538$ clauses
(\texttt{b4d29017\dots} for B, \texttt{1a7c0e70\dots} for A), and $H(1)$ is variable $\binom92+1=37$ of $F_B$ and
$7{,}820+36+1=7{,}857$ of $F_A$. The sevenfold certificate for A is for $F_A$ plus $H(1)$ (\texttt{fa297b5b\dots}, made
by \dirn{C7UPPER/addunits.py}); the one for B is for $F_B$ alone, the case of part~(a$'$). At $(6,8,11)$ the
certificate for B is also for $F_B$ alone, through $64$ cubes, and the one for A is for $F_A$ plus $Y(1,2)$, through
$32$ cubes (Section~\ref{k8s:sec:cube}). Sizes and hashes:
\dirn{C6UPPER/README.md}, \dirn{C7UPPER/README.md}, \dirn{C7UPPER/SHA256SUMS.txt}.

\begin{proof}
Let $S$ be an $(m,K,t)$-admissible set. If $m$ is odd, replace it by the generic, normalised admissible set of
Lemma~\ref{k8s:lem:generic}(b). If $m$ is even, $S$ with its layers is a set of nested regular $m$-gons
(Lemma~\ref{k8s:lem:layer}) in general position, hence generic by Lemma~\ref{k8s:lem:generic}(a); rotate it about $O$ so
that it is normalised, which keeps it admissible with the rotated layers in the same order (Lemma~\ref{k8s:lem:affine})
and keeps it generic.
(a$'$) Give the
variables of $F_B$ their values under $\alpha_S$ (Definition~\ref{k8s:def:assign}); every variable of $F_B$ is a $Y$,
$H$, orientation or chain variable. By reading the source, $F_B$ consists of the header (line 192) and exactly the
clauses appended to \py{cls} at lines 80--81, 137, 141, 147, 175 and 190, that is, the families of
Lemmas~\ref{k8s:lem:order}, \ref{k8s:lem:angB}, \ref{k8s:lem:nestcl} and \ref{k8s:lem:chain}; all are satisfied.
Neither Lemma~\ref{k8s:lem:refl} nor the reflection $\tau$ is used.
(a) By Lemma~\ref{k8s:lem:refl}, $\tau S$ is again a generic, normalised $(m,K,t)$-admissible set, and exactly one of
$\alpha_S$ and $\alpha_{\tau S}$ satisfies the reflection unit; replace $S$ by $\tau S$ if necessary. Then $\alpha_S$
satisfies the reflection unit, and by the argument of~(a$'$) it satisfies $F_B$.
(b) Let $m\in\{6,7\}$, take $S$ as in~(a), so that $\alpha_S$ satisfies the reflection unit, and give the variables of
$F_A$ their values under $\alpha_S$. \py{build} calls \py{fam_order}, \py{fam_angle}, \py{fam_nest} and \py{fam_chain} and nothing else (the
Grassmann--Pl\"ucker family of the docstring is not emitted by the program). The first three are satisfied by
Lemma~\ref{k8s:lem:order} (identical clauses) and Proposition~\ref{k8s:prop:A}, the last by the paragraph on encoder
A's chains.
\end{proof}

\begin{proposition}[encoder A's angular and nesting clauses; written proof (not refereed), new here, and exact
computation]\label{k8s:prop:A}
Let $m\in\{6,7\}$, $4\le K\le2m$, $t\ge3$, and let $S$ be as in Definition~\ref{k8s:def:assign}. Every clause emitted
by \py{fam_angle} and \py{fam_nest} of \dirn{gen.py} on input $(m,K,t)$ is satisfied by $\alpha_S$. Two independent
justifications:
\begin{enumerate}[label=(\roman*)]
\item (written proof (not refereed), with the rounding bound by exact computation; new here) the argument below,
for $m\in\{6,7\}$ and every $K$ and $t$;
\item (exact computation, one program sharing \dirn{gen.py}'s literal translation) only for the two certified instances,
\dirn{C6UPPER/encb/compare_ab2.py} and \dirn{C7UPPER/tests/compare_ab.py} compared, for every triple orbit having a clause
in either encoder and every pattern of its layers, the set of orientation literals forced by A's angular and nesting
clauses with that of B's clauses without chain variables: at $(6,8,11)$ $10{,}417$ clauses each, $5{,}297$ orbits,
$23{,}852$ orbit--pattern pairs, $0$ mismatches (\dirn{C6UPPER/logs/encb_compare.log}); at $(7,8,10)$ $74{,}099$
clauses each, $6{,}221$ orbits, $185{,}549$ pairs, $0$ mismatches (\dirn{C7UPPER/logs/repro_validation.log}). With
Lemmas~\ref{k8s:lem:angB} and~\ref{k8s:lem:nestcl}, an A clause whose premises are true under $\alpha_S$ forces a literal
that a B clause with true premises also forces, hence a true one; an A clause with a false premise is satisfied.
\end{enumerate}
\end{proposition}

\begin{proof}[Proof of (i)]
\py{fam_nest} emits the same nesting units as encoder~B (Lemma~\ref{k8s:lem:nestcl}). In \py{fam_angle}, a pattern is a
permutation \py{perm} of $M_c$ with $0$ first, plus $H$ bits for odd $m$; its premises are the literals
$[\psi_{\mathrm{perm}(r)}<\psi_{\mathrm{perm}(r+1)}]$ for consecutive entries not involving layer $0$ (\py{ylit} makes
those constant) and the $H$ literals. Since $\psi_0=0$ is the least phase, all premises are true under $\alpha_S$ iff
\py{perm} is the $\psi$-order of $M_c$ and the bits are the $h_L$; the emission rule is the same as in
Lemma~\ref{k8s:lem:angB}. For the pattern of $S$, \py{patterns} supplies phases $\psi_L=u\,q_L/(n+1)$, $n=|M_c|$,
$q_L\in\{0,\dots,n\}$ increasing along \py{perm}, and $\theta_L=\psi_L+(1-[\text{bit}])\pi/m$ for odd $m$. These realise the
same $\psi$-order and the same $h$ as $S$ on $M_c$. Lemma~\ref{k8s:lem:keys} applies to both phase assignments on $M_c$
(the sampled $\psi_L$ are distinct and $\theta_L\in[0,2\pi/m)$), and they give the six objects of $c$ the same keys. So the
objects have the same cyclic word, the same counterclockwise order of points and the same block as in $S$. \py{angle_config} sorts the three directions, forms
the gaps, returns \texttt{degenerate} if a gap is within $10^{-9}$ of $0$ or $\pi$, \texttt{in} with the sign of the
sorting permutation if all gaps are $<\pi$, and otherwise the points in order starting after the largest gap; by
Lemma~\ref{k8s:lem:dich} these are exactly B's three cases and outputs, \emph{provided} the floating-point decisions
agree with the exact ones. They do: every exact direction is an integer multiple of $\delta_0=u/(n+1)$, and so is $\pi$
($\pi=m(n+1)\delta_0$ for odd $m$, $\frac{m}{2}(n+1)\delta_0$ for even $m$); every exact direction is at most $2\pi-\delta_0$;
and the three points have distinct directions (distinct $\psi$ for distinct layers, distinct $k$ in one layer). So each
exact gap equals $\pi$ or lies at distance at least $\delta_0\ge\pi/(4m)$ from both $0$ and $\pi$ (exactly, the least
such distance is $\pi/12$ for $(m,t)=(6,11)$ and $\pi/28$ for $(7,10)$). The float inputs depend only on $m$, on
$n\in\{1,2,3\}$, on $q_L$, on the bit and on $k$, not on $K$ or $t$. Re-evaluating \dirn{gen.py}'s float expressions in
all $54$ ($m=6$) and $105$ ($m=7$) such cases and comparing with the exact angles ($\pi$ enclosed by $36$-digit
rationals) shows two things (exact computation, \dirn{k8_soundness_checks.py}). First, each value
\texttt{theta[i] + 2*math.pi*k/m} that \py{angle_config} reduces is a float in $[0,\texttt{2*math.pi})$, so the
reduction \texttt{\%~(2*math.pi)} returns it unchanged (a wrap-around would have produced an error near $2\pi$).
Second, the value used is within $8.7\cdot10^{-16}$ of the exact direction. Hence the sorted order is
the exact one, each computed gap is within $10^{-14}$ of the exact gap (two angle errors, the error of
\texttt{2*math.pi}, and at most two roundings of numbers below $8$, each below $9\cdot10^{-16}$), gaps equal to $\pi$ (only same-layer antipodal
pairs, even $m$) are reported as \texttt{degenerate}, and all other comparisons with \texttt{eps}$=10^{-9}$ and with $\pi$
are decided correctly. The representative's classification is then that of $S$, and the emitted literal is true under
$\alpha_S$ exactly as in Lemma~\ref{k8s:lem:angB}.
\end{proof}

\subsection{The omitted antipodal case}\label{k8s:sec:antipodal}

\begin{lemma}[diameter; written proof (not refereed); the formula is not in the review]\label{k8s:lem:diam}
Let $m$ be even, $q$ and $-q$ points of one layer, and $p$ another point. Then
\[
\orient(q,-q,p)=-2\det(q,p)=-2|q||p|\sin\bigl(\dirv(p)-\dirv(q)\bigr),
\]
so $\orient(q,-q,p)>0$ if and only if $p$ lies strictly clockwise from $q$ by less than $\pi$. For a generic normalised
configuration this is decided by the keys of $q$, $-q$ and $p$ (Lemma~\ref{k8s:lem:keys}), i.e.\ by the phase pattern
alone.
\end{lemma}

\begin{proof}
$\orient(q,-q,p)=\det(-q-q,\,p-q)=-2\det(q,p)+2\det(q,q)=-2\det(q,p)$, and $\det(q,p)=|q||p|\sin(\dirv p-\dirv q)$;
$\dirv(p)\ne\dirv(\pm q)$ by Lemma~\ref{k8s:lem:collinear}.
\end{proof}

A triple containing a same-layer antipodal pair has the centre on one of its sides, so it is neither a triangle around
the centre nor contained in an open half-plane through the centre. For every pattern of such an orbit, encoder~B finds
coinciding keys (line 113) and encoder~A a gap of $\pi$ (\texttt{degenerate}); neither emits an angular clause. For
$(m,t)=(6,11)$ these are $352$ of the $7{,}634$ orbits ($11\cdot3\cdot64$ triples in orbits of size $6$; exact
computation, \dirn{k8_soundness_checks.py}). A nesting triple $\{(i,0),(i,1),(j,k)\}$, $j>i$, contains no
same-layer antipodal pair (as $m\ge4$), and rotations preserve this property, so none of these orbits carries a nesting
unit either; their orientations are constrained only through the chain family. By Lemma~\ref{k8s:lem:diam} valid
clauses for them exist and are omitted. Omitting valid clauses is a
weakening: it can only enlarge the set of satisfying assignments, and Theorem~\ref{k8s:thm:sound} concerns only the
emitted clauses. So the omission is not a source of unsoundness, and the certificates prove unsatisfiability without
these clauses. More generally the encoders omit every orientation fact not covered by rules A and B (for instance
half-plane triples whose middle point is not on the outermost layer, and the two-point-chord facts of
\dirn{CHORD20261105}); each omission is a weakening. The chord-lemma family of \dirn{CHORD20261105} (written proof (not
refereed), each sixfold clause checked by exact computation) contains such diameter clauses; it was not part of the
\dirn{C6UPPER} or \dirn{C7UPPER} formulas. The \dirn{gen.py} docstring states the angular dichotomy without this third
case (finding of \dirn{C6C7REVIEW/SOUNDNESS.md}); the emitted formulas are unaffected.

\subsection{The cube cover and the passage to geometry}\label{k8s:sec:cube}

\begin{lemma}[cube cover; written proof (not refereed); Lean-checked for DIMACS clause lists]\label{k8s:lem:cube}
Let $F$ be a CNF formula and $x_1,\dots,x_n$ variables. If $F\wedge\ell_1\wedge\dots\wedge\ell_n$ is unsatisfiable for each of
the $2^n$ choices $\ell_j\in\{x_j,\neg x_j\}$, then $F$ is unsatisfiable.
\end{lemma}

\begin{proof}
A model $\beta$ of $F$ satisfies the cube with $\ell_j=x_j$ if $\beta(x_j)$ is true and $\ell_j=\neg x_j$ otherwise.
\end{proof}

For formulas given as lists of DIMACS clauses the lemma is the composition of two Lean-checked theorems,
\lean{EncB.not_dimacsSatisfiable_of_cover} applied to \lean{EncB.covers_signPatterns}
(\dirn{ENCBLEAN/EncB/Cover.lean}).
No geometric meaning of the $x_j$ is needed, so sign patterns of the $Y$ variables that contradict transitivity need
no separate argument (their cubes were refuted anyway). Two sixfold splits are used.

\emph{The split of $F_B(6,8,11)$ (encoder~B, $64$ cubes).} Here $F$ is encoder~B's base formula itself, and
$x_1,\dots,x_6$ are $Y(1,2),Y(1,3),Y(1,4),Y(2,3),Y(2,4),Y(3,4)$: variables $1,2,3,10,11,18$ (lines 36--39); $40$ of the
$64$ patterns contradict transitivity. The $32$ cubes with $Y(1,2)$ are encoder~B's cubes of the split below; the $32$
with $\neg Y(1,2)$ were written on 6 November 2026 in the same layout with the first unit negated. That each of the
$64$ cube files is the base formula byte for byte, with adjusted header, plus six unit clauses on these variables, and
that the $64$ patterns are distinct, was checked by \dirn{verify_cover64.py} and, from the files alone, by
\dirn{verify_final64.py} (exact computation; \dirn{C6UPPER/noty12-2026-11-06}). By Lemma~\ref{k8s:lem:cube},
unsatisfiability of the $64$ cube formulas gives that $F_B(6,8,11)$ is unsatisfiable, the first hypothesis of
Corollary~\ref{k8s:cor:geom}(a). For this split no variable needs a geometric meaning. For a Lean definition of
$F_B(6,8,11)$ this application of the lemma is Lean-checked: \lean{EncB.encB_unsat_of_cubes64}, where
\lean{EncB.cubes64} lists the $64$ sign patterns of the variables
$1,2,3,10,11,18$ in the order of the cube files. That the printed texts of that definition and of its $64$ cubes are
the certified files is checked by hashes only (gap~(G2)).

\emph{The split of the base formula plus $Y(1,2)$ (both encoders, $32$ cubes each).} Here $F$ is the
base formula plus the unit $Y(1,2)$, and $x_1,\dots,x_5$ are $Y(1,3),Y(1,4),Y(2,3),Y(2,4),Y(3,4)$: variables
$2,3,10,11,18$ in encoder B and $7{,}636,7{,}637,7{,}644,7{,}645,7{,}652$ in encoder A (by the variable numbering
described in the paragraph after Definition~\ref{k8s:def:assign}); $20$ of the $32$ patterns contradict transitivity.
That each of the $32$ cube files of each encoder is the base formula byte for byte,
with adjusted header, plus $Y(1,2)$ and the five literals of its pattern, and that the $32$ patterns are distinct, was
checked by \dirn{C6UPPER/tools/verify_cover.py} and by the assertion-independent \dirn{C6C7REVIEW/strict_cover.py}
(exact computation). Only the variable of the unit $Y(1,2)$ must have its geometric meaning; it is fixed by the source
(Theorem~\ref{k8s:thm:sound}). This is the route of encoder~A, the cross-check; for encoder~B these $32$ cubes are one
half of the split above.

\begin{corollary}[from unsatisfiability to the Lean propositions; written proof (not refereed), not in the review,
with the cited Lean-checked reduction]\label{k8s:cor:geom}
Let $F_A$, $F_B$ be as in Theorem~\ref{k8s:thm:sound}.
\begin{enumerate}[label=(\alph*)]
\item If $F_B(6,8,11)$ is unsatisfiable, then \lean{Sixfold66} holds; with \lean{sixfold_octagon} (Lean-checked), every
sixfold rotationally symmetric set of at least $61$ points in general position contains a convex octagon. The same
conclusion follows if, for one encoder, the $32$ formulas ``base $(6,8,11)$ + $Y(1,2)$ + cube'' are all unsatisfiable
(this route uses Lemma~\ref{k8s:lem:refl}, through Theorem~\ref{k8s:thm:sound}(a), resp.~(b); it is the route of
encoder~A, the cross-check).
\item If $F_A(7,8,10)$ plus $H(1)$, or $F_B(7,8,10)$, is unsatisfiable, then \lean{Sevenfold70} holds; with
\lean{sevenfold_octagon} (Lean-checked), every sevenfold rotationally symmetric set of at least $65$ points in general
position contains a convex octagon.
\end{enumerate}
\end{corollary}

\begin{proof}
(a) Suppose \lean{Sixfold66} fails: some $T$ in general position, invariant under the rotation by $2\pi/6$ about a
point $c\notin T$, with $|T|=66$, has no convex octagon. Let $\eta(x)=x-c$ be the translation and $R_c$ the rotation in \lean{Sixfold66}
(\lean{rotAbout} with centre $c$ and angle $2\pi/6$), so that $R_c(x)=c+\rho(x-c)$ (with $m=6$). Then $T-c=\eta(T)$ is $(6,8,11)$-admissible: it is in
general position and $8$-free by the first part of Lemma~\ref{k8s:lem:affine}, which holds for every affine bijection
(its last clause assumes $g(O)=O$ and is not used here, since $\eta(O)=-c$); it is $\rho$-invariant because
$\eta R_c\eta^{-1}(y)=R_c(y+c)-c=\rho(y)$, so $\rho(\eta T)=\eta R_c(T)=\eta T$ (here $R_cT=T$, as $R_c$ maps the finite
set $T$ injectively into itself); $O\notin\eta(T)$ because $c\notin T$; and $|\eta(T)|=66$. By
Theorem~\ref{k8s:thm:sound}(a$'$), $F_B(6,8,11)$ is satisfiable, which contradicts the first hypothesis; this route does
not use Lemma~\ref{k8s:lem:refl}. Under the second hypothesis, by Theorem~\ref{k8s:thm:sound}(a), resp.~(b), the base
formula of the encoder plus $Y(1,2)$ is satisfiable; by Lemma~\ref{k8s:lem:cube}, so is one of the $32$ cube formulas,
a contradiction. (b) Replace $2\pi/6$, $66$ and $(6,8,11)$ by $2\pi/7$, $70$ and $(7,8,10)$ in the first part of the
argument; then $T-c$ is $(7,8,10)$-admissible. No cube step is needed: unsatisfiability of $F_A$ plus $H(1)$
contradicts Theorem~\ref{k8s:thm:sound}(b), and unsatisfiability of $F_B$ contradicts
Theorem~\ref{k8s:thm:sound}(a$'$).
\end{proof}

For the Lean definition \lean{EncB.encB} of $F_B$ (gap~(G1)), the first statement of~(a) is Lean-checked without the
written proofs of this section: \lean{EncB.Geo.sixfold66_of_unsat} derives \lean{Sixfold66} from the unsatisfiability
of \lean{EncB.encB} at $(6,8,11)$, through a polar labelling of the peeled layers
(\lean{EncB.Geo.labelling_of_counterexample}), the model of the formula that it defines
(\lean{EncB.Geo.modelOfLabelling}) and \lean{EncB.encB_satisfiable}. With the cube cover, the reading of the printed
texts and \lean{sixfold_octagon} it is part of \lean{EncB.Geo.sixfold_es8_of_texts} (gap~(G2)).

The hypotheses of the corollary are supplied by certificates (cake\_lpr, \texttt{s VERIFIED UNSAT}). Their status,
including how many verdicts are bound to file hashes by a replay against freshly generated formulas, is stated in Section~\ref{sec:cert} and is not repeated here. The link
from the certified files to the programs analysed here is that fresh runs reproduce the formula files byte for byte\footnote{Note added in this paper: for the $64$ sixfold cube files of the hash-bound replay of 6 November 2026 this is on record by hash for all of them, and structurally for all; the $32$ cube files of encoder~B with $\neg Y(1,2)$, written on 6 November, are the regenerated base formula plus six units, checked byte for byte; see Section~\ref{sec:enc}, ``Validation''.}
(hashes above; \dirn{C6UPPER/logs/regen.log}, \dirn{C6C7REVIEW/REPLAY.md}).

\subsection{Trust base and gaps}\label{k8s:sec:trust}

\begin{table}[ht]
\centering\footnotesize
\begin{tabular}{@{}>{\raggedright\arraybackslash}p{0.47\linewidth}>{\raggedright\arraybackslash}p{0.47\linewidth}@{}}
\toprule
Ingredient & Status\\
\midrule
Strict nesting, radius order, centre exclusion (Lemmas~\ref{k8s:lem:centre}, \ref{k8s:lem:nest}); reduction to
\lean{Sixfold66}, \lean{Sevenfold70} & Lean-checked (\dirn{LAYERLEAN})\\
Layer lemma (Lemma~\ref{k8s:lem:layer}) & Lean-checked core; the assembly of the peeled layers is a written proof
(not refereed), not in the review; the part that the Lean route for encoder~B (last row) uses, a polar labelling of
the layers for even $m$, is Lean-checked\\
Signs of rules A and B (Lemma~\ref{k8s:lem:dich}(b),(c)); chain soundness (Lemma~\ref{k8s:lem:chain}, geometry) &
Lean-checked\\
Collinearity with the centre (Lemma~\ref{k8s:lem:collinear}) & written proof (not refereed), new here\\
Convex position via orientations; affine images (Lemmas~\ref{k8s:lem:convpos}, \ref{k8s:lem:affine}; for the case
$g\rho g^{-1}=\rho^{-1}$ see the last row) & written proof (not refereed), standard; not in the review\\
Genericity and perturbation (Lemma~\ref{k8s:lem:generic}) & written proof (not refereed); proves in full the review's
hyperplane statement\\
Orbit variables and stabilisers; key order; circular dichotomy (a); premises (Lemmas~\ref{k8s:lem:orbit},
\ref{k8s:lem:keys}, \ref{k8s:lem:dich}, \ref{k8s:lem:prem}) & written proof (not refereed)\\
Code-to-clause correspondence for both encoders (Lemmas~\ref{k8s:lem:angB}--\ref{k8s:lem:chain},
Theorem~\ref{k8s:thm:sound}) & written proof (not refereed) by reading the source\\
Cube cover (Lemma~\ref{k8s:lem:cube}; $64$ cubes for $F_B(6,8,11)$, $32$ for $F_A(6,8,11)$ plus $Y(1,2)$) & written
proof (not refereed); Lean-checked for DIMACS clause lists (\dirn{ENCBLEAN})\\
Encoder A's float classification, $m\in\{6,7\}$ only (Proposition~\ref{k8s:prop:A}) & exact computation (comparison
with B at the two instances; rounding bound), and written proof (not refereed), new here\\
Cube files equal base plus units; orbit and pattern counts; formula hashes & exact computation\\
Omitted antipodal clauses (Lemma~\ref{k8s:lem:diam}) & written proof (not refereed), formula not in the review; a
weakening only\\
Passage to \lean{Sixfold66}, \lean{Sevenfold70} (Corollary~\ref{k8s:cor:geom}) & written proof (not refereed), not in
the review; for encoder~B at $(6,8,11)$ also Lean-checked end to end, assuming the $64$ cake\_lpr verdicts and
(E1)--(E6) (last row)\\
Unsatisfiability of the formulas & certificate (Section~\ref{sec:cert})\\
\midrule
Not needed for the upper bounds: nested $m$-gons are admissible (Remark~\ref{k8s:rem:ident}) & written proof (not
refereed), new here\\
Used only by encoder~A (cross-check), not for the bounds via encoder~B: the reflection unit
(Lemma~\ref{k8s:lem:refl}) and the case $g\rho g^{-1}=\rho^{-1}$ of Lemma~\ref{k8s:lem:affine} & written proof (not
refereed); Lemma~\ref{k8s:lem:affine} not in the review\\
Lean route for encoder~B at $m=6$ (\dirn{ENCBLEAN}; gaps~(G1), (G2)): for a Lean definition of $F_B$,
satisfaction by the assignment of every combinatorial model (even $m\ge2$), a model from every admissible set (even
$m$; polar labelling of the layers), the $64$-cube cover at $(6,8,11)$, the round trip of its printed text, and the
assembly with \lean{sixfold_octagon} (\lean{EncB.Geo.sixfold_es8_of_texts}) & Lean-checked end to end
(\dirn{ENCBLEAN/EncB}: standard library only; \dirn{ENCBLEAN/geo}: Mathlib \texttt{f897ebcf}), assuming the $64$
cake\_lpr verdicts and the trust items (E1)--(E6) of gap~(G2): that the printed texts are the certified files is
evidence only (compiled code, hashes), and cake\_lpr's parser is tested, not verified; that the definition equals
\dirn{gen_b2.py}'s output is evidence only and not used on this route; not for $m=7$\\
\bottomrule
\end{tabular}
\caption{What the implication ``certificates $\Rightarrow$ \lean{Sixfold66}, \lean{Sevenfold70}'' rests on. The last three
rows are not used by the written proof for encoder~B: the first of them is used by neither encoder, the second only by
encoder~A, and the third is a second route for encoder~B at $m=6$, Lean-checked end to end (assuming the $64$
cake\_lpr verdicts and the trust items (E1)--(E6)) and independent of the written proof.}
\label{k8s:tab:trust}
\end{table}

\paragraph{Gaps that remain.}
\begin{enumerate}[label=(G\arabic*)]
\item The correspondence between the Python sources and the clause families is established by reading the code at the
stated hashes; there is no formal semantics of the programs, and the Python interpreter is trusted to run them as
written. For encoder~B with $m=6$ this gap is closed in Lean on a second route (\dirn{ENCBLEAN}), for a Lean
definition of its formula, \lean{EncB.encB}: a Lean-checked proof that for even $m\ge2$ it is satisfied by the
assignment derived from every \lean{EncB.Model} (\lean{EncB.encB_satisfiable}), and a Lean-checked proof that for even
$m$ every admissible set yields an \lean{EncB.Model}, through a polar labelling of its layers
(\lean{EncB.Geo.labelling_of_counterexample}, \lean{EncB.Geo.modelOfLabelling}; with Mathlib). An \lean{EncB.Model} is a combinatorial interface: an orientation of
each triple of labels and an order of the layer phases, together with, as hypotheses, the properties that the clause
families use (the orientation is alternating and invariant under~$\rho$, the phase order is a strict total order with
layer~$0$ first, rules~A and~B hold on the circular order of the points and their antipodes, the nesting relations
hold, and there is no chain of $K-1$ points around a pivot). On that route the Python source is not used: the Lean
definition is tied to the certified files by the hashes of its printed texts (trust item~(E1) under~(G2)). That
\lean{EncB.encB} equals the output of
\dirn{gen_b2.py} is not proved; it was checked by running compiled code (equal bytes on $222$ of the $225$ instances
of a small grid, a refusal of the Lean program on the other $3$, where \dirn{gen_b2.py} raises, and equal bytes on the
recorded instances, among them $F_B(6,8,11)$). So (G1) stands for the written proof, which argues about
\dirn{gen_b2.py}, and in particular for $m=7$ and for encoder~A; for $m=6$ via encoder~B the Lean route replaces it by
the trust items of~(G2).
\item The passage from a DIMACS file to \lean{Sixfold66} or \lean{Sevenfold70} is a written proof
(Corollary~\ref{k8s:cor:geom}). For encoder~B at $(6,8,11)$ it is formalized end to end for the Lean definition
of~(G1), up to the trust items (E1)--(E6) below (\dirn{ENCBLEAN}): \lean{EncB.Geo.sixfold_es8_of_texts} states that if each of the $64$ Lean strings that
the package's printer function computes for that formula with one of the $64$ cubes is parsed, by a DIMACS parser
written in the package, to an unsatisfiable clause list, then every set of at least $61$ points in general position
that is invariant under the rotation by $2\pi/6$ about some point contains a convex octagon
(\lean{EncB.Geo.sixfold_es8_order6_of_texts}: the same for the rotation by $2\pi/6$ or by $2\pi\cdot5/6$). It
composes the round trip of the printed texts (\lean{EncB.parseDimacs_render}, \lean{EncB.parseDimacs_cube64_get},
\lean{EncB.cubesUnsat64_of_texts_get}), the $64$-cube cover (\lean{EncB.encB_unsat_of_cubes64}), the two Lean
proofs of~(G1) and \lean{sixfold_octagon}. Its hypothesis is equivalent to the unsatisfiability of the $64$ cube
formulas of the Lean definition (\lean{EncB.Geo.texts_iff_cubesUnsat64}), so a failing parse cannot make it hold
trivially; it says exactly what the $64$ verdicts are taken to certify. Not Lean-checked,
and trusted on this route:
\begin{enumerate}[label=(E\arabic*)]
\item that these $64$ texts are the certified cube files: the SHA-256 of the output of the compiled printer equals the
record for the base formula and for all $64$ cube files (\dirn{ENCBLEAN/scripts/hash_binding.sh},
\dirn{ENCBLEAN/logs/m9_a2fix_hash_binding.log}); evidence from running compiled code, not a Lean theorem;
\item the $64$ hash-bound cake\_lpr verdicts \texttt{s VERIFIED UNSAT} on these files, and the provenance of the
cake\_lpr binary that gave them (the proof files are not public); the records that tie each verdict to the hash of
its file were written by the Python harnesses of the two runs (\dirn{C6UPPER/replay-2026-11-06/replay.py},
\dirn{C6UPPER/noty12-2026-11-06/noty12.py}) and cross-checked by \dirn{verify_final64.py};
\item that cake\_lpr reads each of these files as the package's parser does, and that its \texttt{s VERIFIED UNSAT}
means unsatisfiability of the clause list it read; the second half rests on cake\_lpr's verified HOL4/CakeML
checker, the first is not verified (tested on samples only, \dirn{ENCBLEAN/logs/m9_a2_cake_parse.log});
\item the Lean compiler and runtime, for the printing program of~(E1) (not for the theorems), and a reading of that
program's input/output code (that it prints exactly these texts and nothing else);
\item the Lean kernel, Mathlib \texttt{f897ebcf}, and the match of the Lean definitions (\lean{GeneralPosition},
\lean{HasConvexNGon}, \lean{orient}, \lean{rotAbout}, \lean{Invariant}) with the notions of
Section~\ref{k8s:sec:setting};
\item SHA-256.
\end{enumerate}
So the sixfold bound is formally verified end to end via encoder~B, assuming the $64$ cake\_lpr verdicts and the trust
items (E1)--(E6). The results of this section other than the cited Lean theorems are written proofs only, and the
sevenfold bound and the route of encoder~A are not formally verified end to end.
\item Encoder A's angular part rests on binary64 arithmetic. For the two certified instances
Proposition~\ref{k8s:prop:A} covers it twice; for other $(K,t)$ with $m\in\{6,7\}$ only argument~(i) applies, and for
other $m$ no claim about encoder~A is made. Argument~(i) is new here and unreviewed, and comparison~(ii) is a single
program that uses \dirn{gen.py}'s own literal translation.
\item The written proofs marked \emph{new here} (Lemma~\ref{k8s:lem:collinear}, Remark~\ref{k8s:rem:ident} and
Proposition~\ref{k8s:prop:A}(i)) and those marked \emph{not in the review} (Lemmas~\ref{k8s:lem:convpos}
and~\ref{k8s:lem:affine}, the assembly in Lemma~\ref{k8s:lem:layer}, the formula of Lemma~\ref{k8s:lem:diam}, and
Corollary~\ref{k8s:cor:geom}) have had no internal review. The remaining written proofs expand arguments reviewed on
4 November. Part~(a$'$) of Theorem~\ref{k8s:thm:sound} and the $64$-cube split of Section~\ref{k8s:sec:cube} were added
on 6 November, after that review; the proof of~(a$'$) is that of~(a) without the reflection step.
Lemma~\ref{k8s:lem:refl} and the case $g\rho g^{-1}=\rho^{-1}$ of Lemma~\ref{k8s:lem:affine} are used only by
encoder~A (cross-check); the bounds via encoder~B (Theorem~\ref{k8s:thm:sound}(a$'$), Corollary~\ref{k8s:cor:geom})
do not depend on them. The Lean route for encoder~B at $m=6$ of~(G2) uses none of the written proofs of this section.
\end{enumerate}

\paragraph{Where the sources only sketched.} The docstring of \dirn{gen.py} (step~0) calls $K$-freeness a strict
inequality in the parameters; Lemma~\ref{k8s:lem:convpos} supplies the needed fact that convex position is decided by
orientation signs. The review states that the bad phase equalities form a finite union of proper affine hyperplanes
near the configuration; Lemma~\ref{k8s:lem:generic}(b) proves this in full and identifies the perturbed set as an
admissible set through Lemma~\ref{k8s:lem:layer}. The angular dichotomy is completed by the antipodal case
(Lemma~\ref{k8s:lem:keys} and Section~\ref{k8s:sec:antipodal}), as the review asked. The chain argument of the docstring
and of the review (a simple, locally convex fan polygon) is replaced by the Lean-checked proof through positive
triples. Stating everything for admissible sets with their peeled layers lets the Lean theorems apply to the layers
directly; the link to nested $m$-gons given abstractly is Remark~\ref{k8s:rem:ident}, which the upper bounds do not
need.
